\documentclass[a4paper]{article}
\usepackage{luatexja}
\usepackage[left=10truemm,right=10truemm,top=25truemm,bottom=20truemm]{geometry}
\usepackage{mathtools}
\usepackage{amsmath}
\usepackage{amsfonts}
\usepackage{bm}
\usepackage{setspace}
\usepackage{wrapfig}
\usepackage{docmute}
\usepackage{amssymb}
\usepackage{hyperref}
\DeclareMathOperator\arctanh{arctanh}
\DeclareMathOperator\arccosh{arccosh}
\newcommand{\Jnorm}{J_{\rm norm}}
\newcommand{\Jbound}{\frac{3\pi^2}{8}}

\title{二重時空理論：\\通常時空と双対時空の間のねじれとしての重力}

\author{https://github.com/hypernumbernet}
\date{\today}

\begin{document}

\maketitle

\begin{abstract}
一般相対性理論（GR）は重力を単一の時空連続体の曲率として記述する。その経験的成功にもかかわらず、量子力学や特異点の本性と対峙すると深い困難に直面する。本論文では根本的な再定式化を提案する。重力は連続多様体の曲率ではなく、すべての質量粒子に内在的に付随する2つのコンパクト時空の相対的なずれから生じるねじれである。

16実次元の双四元数代数（$\mathrm{Cl}(3,1)$ と同型）を用いて、基底 $(j, kI, kJ, kK)$ による通常時空と基底 $(k, jI, jJ, jK)$ による双対時空を定義する。双対写像 $X \mapsto Xi$ は時間の矢を反転させつつミンコフスキーノルム $X^2 = X'^2$ を保つ。ローレンツ変換および局所フレーム回転は、各時空に別々に作用する可換なローターの対によって生成される。

完全ローターは $R_\text{total} = \exp[(\phi_1/2)iI + (\phi_2/2)iJ + (\phi_3/2)iK + (\theta_1/2)I + (\theta_2/2)J + (\theta_3/2)K]$ である。
ねじれ不整合ローター $\Omega = R^\dagger_\text{usual} R_\text{dual}$ からねじれバイベクトル $\Omega_\text{biv} = \log \Omega$ が得られる。
本論文を通じて用いる一意のローレンツ不変スカラーは、6生成子のローレンツ候補 $\mathfrak{so}(3,1)$（12次元直和 $\mathfrak{so}(3,1)\oplus\mathfrak{so}(3,1)$ ではない）上のパラメータ空間形式 $J=\frac12\sum_a(\alpha_a^2-\beta_a^2)$ である。半角展開 $\Omega_{\mathrm{biv}}=\sum_a\bigl(\frac{\alpha_a}{2}i\Gamma_a+\frac{\beta_a}{2}\Gamma_a\bigr)$ の下で $B(\Omega_{\mathrm{biv}},\Omega_{\mathrm{biv}})=2\sum_a(\alpha_a^2-\beta_a^2)$ となり、したがって $J$ は $\frac1{16}B(\Omega_{\mathrm{biv}},\Omega_{\mathrm{biv}})$ の4倍である。許容配置においては $|J|\le\Jbound$、あるいは同値に $\Jnorm=\frac8{3\pi^2}J$ として $|\Jnorm|\le 1$ である。

作用 $S = \frac{c^4}{16\pi G} \int J \, d^4x$ はアインシュタイン--ヒルベルト作用と動学的に等価であり、GR を正確に再現しつつクリストッフェル記号、リーマン曲率、および連続体時空仮説そのものを排する。本理論は一般相対性理論の遠平行等価理論（TEGR）の双四元数的実現であり、ねじれは今や双対時空ねじれとして物理的に解釈される。

重要なのは、$J$ が粒子内在的な双対ローター間の相対角のみに依存するため、双対時空成分 $\theta_a$ のコヒーレント励起が重力工学への直接的経路を開くことである。重力遮蔽、慣性質量の低減、さらには反重力も、初めて基本法則の違反ではなくローター同期の問題となる。

GR は置き換えられない。完成され、重力は原理上制御可能になる。
\end{abstract}

\section{はじめに}
\label{sec:introduction}

1915年の誕生以来、一般相対性理論は物理学における重力記述の最も成功した理論の一つとして立っている。しかし水星の近日点進行から重力波・ブラックホール撮像に至る経験的成功にもかかわらず、深い哲学的な不安が残る。なぜ物質は時空を曲げ、なぜ曲がった時空が物質の運動を支配するのか。アインシュタイン場の方程式は重力がどのように働くかを極めて精密に与えるが、曲率そのものの起源についてより深い理由は与えない。この説明のギャップは、マッハの原理と対峙すると特に鋭くなる。物体の慣性は宇宙の物質の総体によって決まるべきであるのに対し、GR は時空を完全な空虚さでも存在しうる独立した連続体として扱う。

ループ量子重力、弦理論、創発的重力など、これらの問題を解決しようとする一世紀の試みは数学的優雅さをもたらしたが、時空の微視的起源や質量--エネルギーが曲率を生成する物理機構についての合意には至っていない。いずれのアプローチも連続体仮説を保持する。時空は先験的に滑らかで微分可能な多様体であり、重力はクリストッフェル記号とリーマンテンソルを通じて計量の平坦さからのずれとして実現される。

本稿はこの伝統と断固として決別する。

時空連続体仮説は根本的に誤りであると提案する。共有される単一の時空多様体は存在しない。代わりに、各粒子が独自の二重時空の対を内在的に運ぶ。それは双四元数代数（クリフォード代数 $\mathrm{Cl}(3,1)$ と同型）に符号化された16実次元のコンパクト化された構造である。通常時空（基底 $j, kI, kJ, kK$）とその双対（基底 $k, jI, jJ, jK$）は右乗算 $i$ で結ばれ、この変換はミンコフスキーノルムを保ちながら時間の矢を反転させる。重力と慣性は背景多様体の曲率からではなく、異なる粒子間で平行移動を要求するときの、これら2つの粒子内在的時空のねじれ不整合角から生じる。

この哲学的転換は深い。標準的 GR において時空は物質に運動を教える自律的存在であり、物質は時空に曲がり方を教える（ウィーラーの有名な言葉）。二重時空理論には自律的時空は全くない。時空の自由度はスピンや電荷のように粒子自身の内部属性へと格下げされ、巨視的世界の見かけの曲率は、隣接粒子の双対ローター間のねじれ不整合の集合効果としてのみ現れる。アインシュタイン方程式は正確に回復されるが、クリストッフェル記号、共変微分、リーマンテンソルは完全に消える。それらは双四元数上の単純な代数演算（ローターの指数・対数・キリング形式のトレース）に置き換えられる。

したがって本理論は同時に三つの革命的帰結を達成する。

1. GR との観測上の完全な等価性：ニュートン極限、ブラックホール、宇宙論、重力波を含むすべての予測において、従来の意味での幾何曲率はなくなる。

2. 重力と慣性の自然な統一：等価原理は自由落下において通常・双対ローターが並行であることの要求として現れ、加速度はこの並行性を破り、重力と同一のねじれ現象として慣性力を生む。

3. 重力の原理上の制御可能性：不変スカラー $J$ が粒子内在的な双対ローター間の相対角のみに依存するため、双対時空成分 $\theta_a$ をコヒーレントに励起できる外場は重力相互作用を減少、打ち消し、さらには反転させうる。重力は侵すべからざる基本力ではなく工学的自由度となる。

これは GR の修正ではなく完成である。連続体は有用なフィクションであった。美しいが最終的には不要である。時空が粒子局所的な対構造として認識されれば、「なぜ物質は時空を曲げるのか」という一世紀の謎は解ける。物質は時空を曲げない。自身の二重時空をずらし、そのずれが我々が重力と知覚するものである。

本論文の構成は以下のとおりである。

第\ref{sec:math}節では、$\mathrm{Cl}(3,1)$ と同型の双四元数代数を導入し、各粒子内在的な双対時空の対を構成して数学的基礎を与える。

第\ref{sec:kinematics}節では、行列やクリストッフェル記号なしにローレンツブースト、回転、平行移動を統一する完全ローター形式を展開する。

第\ref{sec:gravity}節では、相対ローター $\Omega = R_\text{usual}^\dagger R_{\text{dual}}$ のみから構成されるねじれ不整合スカラー $J$ の作用がアインシュタイン--ヒルベルト作用と動学的に等価であることを示し、理論が遠平行重力の双四元数的実現であることを明らかにする。

第\ref{sec:unified}節では、慣性と重力が同一の双対ローター剛性から生じる同一現象であり、等価原理が代数的恒等式として現れることを示す。

第\ref{sec:strongfield}節では、恒星残骸を引力・斥力の多層構造をもつねじれ星として再解釈し、特異点と事象地平の両方を禁止する。

第\ref{sec:timedilation}節では、ねじれ不整合から時間膨張効果を導き、強重力場および惑星内部で GR からの測定可能なずれを予測する。

第\ref{sec:nuclear}節では、ねじれ層のスケール不変性を証明し、核力を層をもつ重力として同定し、原子核を原始ねじれ星とみなす。

第\ref{sec:electronshells}節では、電子殻を電気反転層としてモデル化し、微視的スケールの電磁気学の観点から電子軌道と原子間力を説明する。

第\ref{sec:unification}節では、双四元数代数内のすべての基本相互作用の統一の骨子を示し、ゲージ場が付加的なローター成分として現れることを概説する。

第\ref{sec:darkmatter}節では、重力チャンネルと電磁チャンネルの統一を踏まえ、銀河スケールの双対時空コンパクト性とコンパクトから平坦への写像の幾何歪みから平坦な回転曲線を導き、暗黒物質の必要性を排除する。

第\ref{sec:gravitycontrol}節では、円偏波電磁場による双対時空ローターのコヒーレント励起を通じた重力工学的機構を詳述する。

第\ref{sec:baryonasymmetry}節では、初期宇宙における双対時空ねじれ動学に基づくバリオン非対称の機構を提案する。

第\ref{sec:discretetime}節では、双対セクター角が量子化される双対時空の離散モデルを導入する。$J$ の有界性は許容錐上でのみ成り立ち、有限であるのはトーラスの離散ローター像であって、整数環の単数群ではない。

第\ref{sec:dualrotordynamics}節では、外場下での双対ローターの動力学を展開し、ねじれ構成における運動方程式と安定性基準を導く。

第\ref{sec:quantumgravity}節では、理論を量子重力へ拡張し、量子状態を双対時空ローター位相と同定し、ねじれコヒーレンス条件から不確定性関係を導く。

第\ref{sec:conclusions}節では、16次元双四元数代数内のすべての力の統一を要約し、今後10年以内に理論を確認または反証しうる近接実験を概説する。

この再定式化により、GR は連続体の牢獄から解放され、重力が初めて原理上制御可能な理論へと高められる。

\section{数学的基礎：16次元双四元数代数}
\label{sec:math}

二重時空理論全体は、ミンコフスキー時空（符号 $(-,+,+,+)$）上のクリフォード幾何代数 $\mathrm{Cl}(3,1)$ と標準的に同型である双四元数（二重四元数とも呼ばれる）の16実次元代数の中に構築される。この代数は、すべての質量粒子に内在的に付随する単一代数構造の中に、通常と双対という二つの可換なミンコフスキー時空の完全なコピーを自然に収容する。

\subsection{双四元数代数}

まずハミルトンの四元数代数の二つの独立したコピーから出発する。
\begin{itemize}
  \item 生成子 $i,j,k$ が $i^2=j^2=k^2=-1$, \ $ij=k$, \ $ji=-k$（巡回）を満たす主四元数、
  \item 生成子 $I,J,K$ が同一の規則 $I^2=J^2=K^2=-1$, \ $IJ=K$, \ $JI=-K$（巡回）を満たす従四元数。
\end{itemize}
完全な双四元数代数はテンソル積 $\mathbb{H} \otimes_{\mathbb{R}} \mathbb{H}$ であり、二つの集合は厳密に可換である。
\[
iI = Ii,\; iJ = Ji,\; iK = Ki,\; jI = Ij,\; \dots
\]
得られる16個の実線形独立基底は
\[
1,\; i,\; j,\; k,\; I,\; J,\; K,\; iI,\; iJ,\; iK,\; jI,\; jJ,\; jK,\; kI,\; kJ,\; kK.
\]

\subsection{Cl(3,1) との明示的同型}

次の1次（ベクトル）生成子の同定により、代数は $\mathrm{Cl}(3,1)$ と同型である。
\begin{align*}
  e_0 &= j,                    & &\text{（時間様基底ベクトル）}\\[4pt]
  e_1 &= kI, \quad e_2 = kJ, \quad e_3 = kK & &\text{（空間様基底ベクトル）},
\end{align*}
これは直ちに定義関係
\[
e_0^2 = -1, \quad e_1^2 = e_2^2 = e_3^2 = +1, \quad \{e_\mu,e_\nu\} = 0 \;\; (\mu \neq \nu).
\]
を満たす。高次の元は標準的なクリフォード積に従う。特に：
\begin{alignat*}{3}
  &\text{2次（バイベクトル）:}  &\quad& e_0e_1 = iI, \;\; e_0e_2 = iJ, \;\; e_0e_3 = iK, \\
  &&& e_3e_2 = I,   \;\; e_1e_3 = J,   \;\; e_2e_1 = K, \\[6pt]
  &\text{3次（トライベクトル）:} && e_1e_2e_3 = k, \;\; e_0e_3e_2 = jI, \;\; e_0e_1e_3 = jJ, \;\; e_0e_2e_1 = jK, \\[6pt]
  &\text{擬スカラー:}        && e_0e_1e_2e_3 = i.
\end{alignat*}
右乗算 $i$ による時間反転の下で補完的となる「双対」基底は
\[
\tilde{e}_0 = k = e_1e_2e_3, \quad \tilde{e}_1 = jI = e_0e_3e_2, \quad \tilde{e}_2 = jJ = e_0e_1e_3, \quad \tilde{e}_3 = jK = e_0e_2e_1,
\]
で生成され、以下で導入する双対時空表示を与える。

この明示的同型により、16次元双四元数代数はアドホックな構成ではなく、4次元ミンコフスキー時空の幾何的に自然なクリフォード代数であることが明らかになる。それは通常セクターと双対セクターを区別しつつ完全なローレンツ構造を保つ内在的デュアリティで豊かになる。

クリフォード代数形の見かけの不規則性。例えば $j$ だけが唯一の時間様ベクトル生成子であることや、空間基底が交差項優勢の $kI,kJ,kK$ であることは偶然ではない。二つの時間反転したミンコフスキー時空のコピーを同時に符号化する構造の中で、3次元物理空間の完全な回転対称性を保つために正確に存在する。この微妙な非対称性が、双対写像 $X \mapsto Xi$ が時間の矢を反転させても空間等方性を破らないことを可能にし、等価原理と、粒子内在的な双対フレーム間の相対ずれとしての重力学的ねじれの創発のための代数的基礎を与える。

\subsection{粒子内在的な双対時空構造}

各粒子は、双四元数代数内に符号化されたコンパクト化されたミンコフスキー時空の対を内在的に運ぶと仮定される。
\begin{itemize}
  \item \textbf{通常時空：} 次の形のベクトルで表される
  \[
  X = ct \, j + x \, kI + y \, kJ + z \, kK,
  \]
  ミンコフスキーノルムは
  \[
  X^2 = -(ct)^2 + x^2 + y^2 + z^2.
  \]

  \item \textbf{双対時空：} 次の形のベクトルで表される
  \[
  X' = ct' \, k + x' \, jI + y' \, jJ + z' \, jK,
  \]
  同一のミンコフスキーノルム
  \[
  X'^2 = -(ct')^2 + x'^2 + y'^2 + z'^2.
  \]
\end{itemize}
二つの時空は双対写像
\[
X' = X i,
\]
で結ばれ、時間の矢を反転させる。
\[
(ct \, j) i = ct \, (j i) = -ct \, k,
\]
ミンコフスキーノルムは保たれる。$\mathrm{Cl}(3,1)$ の擬スカラー $i = e_0 e_1 e_2 e_3$ はすべての1次ベクトルと反可換（$X \in \mathrm{Cl}(3,1)$ に対し $i X = -X i$）であり $i^2 = -1$ を満たすため、ノルムは次のように変換される。
\[
X'^2 = (X i)(X i) = X (i X) i = -X (X i) i = -X^2\, i^2 = -X^2 \cdot (-1) = X^2.
\]
この双対構造はすべての粒子に内在的である。通常時空は標準的運動学を支配し、双対時空は重力と慣性の存在下で動学的に関連する補完的自由度を符号化する。二つの時空の存在により、ローレンツ変換、平行移動、そして重力の本性のより豊かな幾何解釈が可能となる。

\section{運動学的構造：ローター、ブースト、平行移動}
\label{sec:kinematics}

本節では理論の運動学的枠組みを展開する。すべてのローレンツ変換および局所フレーム回転は、通常時空と双対時空に独立に作用する可換なローターの対によって生成される。要点は、従来双曲回転とされるブーストが、平方が$+1$となるバイベクトルから自然に生じ、虚角度のアドホック導入なしにローターの馴染みの $\cosh/\sinh$ 展開を与えることである。

\subsection{通常時空におけるローター}

運動学的構造を示すため、ラピディティ $\phi$ の $x$ 方向ローレンツブーストを考える。生成バイベクトルは $iI$ であり $(iI)^2 = i^2 I^2 = (-1)(-1) = +1$ を満たす。一般に通常時空内のバイベクトル $iI,iJ,iK$ はすべて平方が $+1$ でブースト様の性質を与え、双対時空内の $I,J,K$ は平方が $-1$ で純回転を想起させる。

対応するブーストローターはしたがって
\[
R_\text{usual} = \exp\!\left( \frac{\phi}{2} iI \right).
\]
生成子の正の平方により、指数は双曲関数へと自然に展開される。
\[
R_\text{usual} = \cosh \frac{\phi}{2} + iI \sinh \frac{\phi}{2}.
\]
このローターをサンドイッチ積で時空ベクトル $X = ct \, j + x \, kI + \cdots$ に適用すると、正準ローレンツ変換が得られる。
\begin{align*}
ct' \, j + x' \, kI &= R_\text{usual} \, X \, R_\text{usual}^\dagger \\
&= \left[ \cosh\frac{\phi}{2} + iI \sinh\frac{\phi}{2} \right] (ct \, j + x \, kI) \left[ \cosh\frac{\phi}{2} - iI \sinh\frac{\phi}{2} \right] \\
&= (\gamma ct + \gamma v x) \, j + (\gamma x + \gamma v ct) \, kI,
\end{align*}
ここで $\gamma = \cosh \phi$, $v/c = \tanh \phi$。通常時空内の純回転は $I,J,K$ のような平方 $-1$ のバイベクトルから生じ、馴染みの $\cos,\sin$ 展開を与える。

\subsection{完全ローター}

完全な双四元数構造上の最も一般の固有時順方向保つ変換は完全ローター
\[
R_\text{total} = \exp\left( \frac{\phi_1}{2} iI + \frac{\phi_2}{2} iJ + \frac{\phi_3}{2} iK + \frac{\theta_1}{2} I + \frac{\theta_2}{2} J + \frac{\theta_3}{2} K \right),
\]
によって生成される。ここでパラメーター $\phi_a$ は通常時空のブーストと回転を、$\theta_a$ は双対側のそれを支配する。通常と双対の生成子は可換（例：$[iI, I] = 0$）であるため、ローターはきれいに因数分解される。
\[
R_\text{total} = R_\text{usual} \, R_\text{dual}.
\]

幾何的意味をより明瞭にするため、大きさと方向を分離し明示的な総和の添え字を消すバイベクトル形のローターを導入する。

通常時空の単位ブーストバイベクトルを
\[
\hat{p} = p_1 \, iI + p_2 \, iJ + p_3 \, iK, \qquad \hat{p}^2 = +1,
\]
と定義する。$\hat{p}$ は正規化（$|\hat{p}| = 1$）され、$\phi = \sqrt{\phi_1^2 + \phi_2^2 + \phi_3^2}$ は全ラピディティ（方向は $p_a = \phi_a / \phi$ に符号化）である。

同様に双対時空の単位回転バイベクトルを
\[
\hat{q} = q_1 \, I + q_2 \, J + q_3 \, K, \qquad \hat{q}^2 = -1,
\]
と定義する。$\hat{q}$ は正規化され、$\theta = \sqrt{\theta_1^2 + \theta_2^2 + \theta_3^2}$ は全双対回転角（方向は $q_a = \theta_a / \theta$）である。

完全ローターは次のエレガントなバイベクトル指数形をとる。
\[
R_\text{total} = \exp\!\left( \frac{\phi}{2} \hat{p} + \frac{\theta}{2} \hat{q} \right)
= R_\text{usual} \, R_\text{dual},
\]
ここで
\[
R_\text{usual} = \exp\!\left( \frac{\phi}{2} \hat{p} \right)
= \cosh\frac{\phi}{2} + \hat{p} \sinh\frac{\phi}{2},
\]
\[
R_\text{dual} = \exp\!\left( \frac{\theta}{2} \hat{q} \right)
= \cos\frac{\theta}{2} + \hat{q} \sin\frac{\theta}{2}.
\]

この統一バイベクトル形は、二つの時空上の完全ローレンツ群作用を同時に符号化し、六つの独立実パラメーター、すなわち通常時空のラピディティ $\phi$ と単位ブースト方向 $\hat{p}$、および双対時空の回転角 $\theta$ と単位回転軸 $\hat{q}$ を明示する。

\subsection{サンドイッチ積}

ヴァーソル形式はこれらの変換を適用する簡潔な手段を与える。時空事象を表す任意の双四元数マルチベクトル $X$ に対し、変換後のベクトルは
\[
\tilde{X} = R_\text{total} \, X \, R_\text{total}^\dagger.
\]
この単一の操作はミンコフスキーノルム（$X^2 = \tilde{X}^2$）を保ち、通常成分と双対成分の双方に一貫して作用し、ブーストと回転を一つの代数的表現に継ぎ目なく統合する。行列や座標依存の機構は不要であり、幾何は代数の内に内在的に展開される。

\section{双対時空間のねじれ不整合としての重力}
\label{sec:gravity}

本中心節では、重力が時空の曲率でも外力でもなく、すべての質量粒子が内在的に運ぶ通常時空と双対時空の間のねじれ不整合であることを確立する。アインシュタイン--ヒルベルト作用はこの不整合から代数的に現れ、クリストッフェル記号、リーマンテンソル、連続多様体への言及は一切不要である。得られる理論は TEGR と数学的に等価だが、深い物理的再解釈を与える。ねじれは今や粒子局所的な双対ローター間の相対回転角である。

\subsection{ねじれ不整合ローターの定義}

ねじれ不整合ローターは
\[
\Omega = R_\text{usual}^\dagger R_\text{dual}.
\]
と定義される。

ダガー操作は双四元数代数におけるローターの標準的共役逆である。
\[
R_\text{usual}^\dagger = \exp\!\left( -\frac{\phi}{2} \hat{p} \right)
= \cosh\frac{\phi}{2} - \hat{p} \sinh\frac{\phi}{2}.
\]
形式的には単純な共役だが、このダガーは二重時空枠組みにおいて特別な意味を持つ。$\sinh$ 項の前の負符号は双対写像に内在する時間反転を正確に符号化する。擬スカラー $i$ は4次でありしたがってバイベクトルと可換であるため、右乗算 $i$ はブースト生成子に対して
\[
\hat{p}\, i = - (p_1 I + p_2 J + p_3 K) \equiv -\hat{q}',
\qquad (\hat{p}\, i)^2 = -1,
\]
と作用する。ここで $\hat{q}' \equiv p_1 I + p_2 J + p_3 K$ は双対セクターにおける単位回転生成子である。同値に $\hat{p} = \hat{q}'\, i$ であり、逆通常ローターは厳密な書き換え
\[
R_\text{usual}^\dagger = \cosh\frac{\phi}{2} - \hat{q}'\, i\, \sinh\frac{\phi}{2}.
\]
を許す。標準的恒等式 $\cosh(\phi/2) = \cos(i\phi/2)$ および $\sinh(\phi/2) = \sin(i\phi/2)/i$ を用いると、これは解析接続（$\phi \to i\phi$）された三角関数形
\[
R_\text{usual}^\dagger = \cos\!\left(\frac{i\phi}{2}\right) - \hat{q}'\, \sin\!\left(\frac{i\phi}{2}\right),
\]
を与え、ダガー操作が通常時空の双曲幾何と双対時空の楕円幾何を橋渡しすることを示す。実引数に対する $\cosh$ と $\cos$ の等価ではなく、双対写像 $X \mapsto X i$ を横断するブーストローターの一意な解析接続である。

相対ローター $\Omega = R_\text{usual}^\dagger R_\text{dual}$ はこの橋渡しを操作的に実現する。ダガーは通常時空の双曲ブースト構造を一時的に「分解」し（符号反転を通じて双対写像を誘起し）、続く $R_\text{dual}$ との乗算が結合対象を「再構成」し、相対的ずれを明示する。この「分解--再構成」過程を通じて、通常時空と双対時空のねじれ不整合が代数的に抽出され、この不整合が二重時空枠組みにおける重力の本質をなす。

真空極限（自由落下）では、完全な反同期が $R_\text{dual} = R_\text{usual}^\dagger$（すなわち $\theta = \phi$ かつ $\hat{q} = -\hat{q}'$）を要求し、$\Omega = 1$ およびねじれの消失を与える。質量--エネルギーはこの同期からのずれを誘起し、非自明な $\Omega \in \mathrm{Spin}^+(3,1) \oplus \mathrm{Spin}^+(3,1)$ を生む。

関連するねじれバイベクトルは
\[
\Omega_\text{biv} = \log \Omega \in \mathfrak{so}(3,1),
\]
である。ここで6個の生成子 $\{i\Gamma_a,\Gamma_a\}$ はローレンツ候補を張り、12次元直和 $\mathfrak{so}(3,1)\oplus\mathfrak{so}(3,1)$ ではない。対数は $\mathrm{Spin}^+(3,1)$ の恒等元の単連結近傍における主枝であり、ベーカー--キャンベル--ハウスドルフ（BCH）展開が収束する自然な領域である。本論文で扱う物理的関心領域は弱場／小不整合セクター $\Omega \approx 1$ であり、$\Omega_\text{biv}$ は一意に定まる。この近傍の外にある位相的に異なる枝は、ここで提示する動力学では用いない。

\subsection{双対ねじれからのローレンツ不変スカラー}

$\mathfrak{so}(3,1)$（ベクトル表現）上のキリング形式は
\[
B(X,Y) = 4\,\mathrm{Tr}(X Y).
\]
\S\ref{sec:kinematics} の 2 次空間バイベクトルをまとめて $\Gamma_a$（$a=1,2,3$）と書き、$\Gamma_1 = I$, $\Gamma_2 = J$, $\Gamma_3 = K$ とする（\S\ref{sec:math} のクリフォード同型 $e_3e_2 = I$, $e_1e_3 = J$, $e_2e_1 = K$ に対応）。通常時空のブースト生成子は $i\Gamma_a$、双対セクターの回転を生成するのは $\Gamma_a$ である。キリング形式を評価すると $B(i\Gamma_a, i\Gamma_a) = +8$, $B(\Gamma_a, \Gamma_a) = -8$ となり、ローレンツ符号を完璧に再現する。同軸生成子は可換であり $[i\Gamma_a,\Gamma_a]=0$ である。異軸の交換子は必ずしも消えない。

半角展開 $\Omega_\text{biv}=\sum_a\bigl(\frac{\alpha_a}{2}i\Gamma_a+\frac{\beta_a}{2}\Gamma_a\bigr)$ と書く。このとき
\[
B(\Omega_\text{biv},\Omega_\text{biv})=2\sum_{a=1}^3(\alpha_a^2-\beta_a^2),
\]
であり、初期草稿の因子 $8\sum$ ではない。本論文を通じて用いるパラメータ空間スカラーは
\[
J=\frac12\sum_{a=1}^3(\alpha_a^2-\beta_a^2),
\]
であり、この展開の下で $\frac1{16}B(\Omega_\text{biv},\Omega_\text{biv})$ の4倍である。有界性の主張はこの $J$ および $\Jnorm=\frac8{3\pi^2}J$ に関するものである。許容配置（非負ラピディティかつ $\alpha_a+\beta_a\le\pi/2$）においては $|J|\le\Jbound$ かつ $|\Jnorm|\le 1$ である。主枝条件のみを課した場合、この有界は成り立たない。

真空では $J = 0$、通常の引力重力では $J > 0$（通常時空ブーストが支配的）、双対時空回転が支配するとき $J < 0$（斥力重力）が可能である。

\subsection{アインシュタイン--ヒルベルト作用との正確な等価性}

作用として
\[
S = \frac{c^4}{16\pi G} \int J \, d^4x
\]
を提案する。この作用は純粋に代数的であり背景独立である。

スカラー $J$ は、全発散を除いて、かつローターからテトラッドへの辞書が固定された後にのみ、TEGR のねじれスカラー $T$ と同定される。同定 $J=\tfrac12 T+\mathrm{div}$ はここでは再導出しない。TEGR 文献における標準的結果として
\[
\int T \, e \, d^4x = -\int R \, \sqrt{-g} \, d^4x + \mathrm{boundary},
\]
が成り立つので、真空および物質結合した場の方程式は正確にアインシュタイン方程式である。GR のすべての解（シュヴァルツシルト、カー、FLRW、重力波など）は精密に回復される。

したがって GR は完全に再現されるが、重力は粒子内在的な双対時空の間のねじれ不整合角へと格下げされる。

\subsection{弱場およびニュートン極限}

線形化領域では
\[
R_\text{usual} \approx 1 + \frac{1}{2} \sum \alpha_a \, i\Gamma_a, \quad R_\text{dual} \approx 1,
\]
となり、$\Omega_\text{biv} \approx \sum \alpha_a \, i\Gamma_a$ および $J \approx \frac{1}{2} \sum \alpha_a^2$ を与える。標準的な TEGR 線形化は $\alpha_a \propto \partial_a (\Phi/c^2)$ を示し、正しい係数をもつポアソン方程式
\[
\nabla^2 \Phi = 4\pi G \rho
\]
に至る。ポストニュートンパラメーターは GR と同一である（$\gamma = \beta = 1$ など）。

\subsection{双対時空におけるローター作用}
\label{subsec:dual-rotors}

二重時空理論において、唯一の物理的実在は、すべての質量粒子が運ぶ内在的双四元数代数 $\mathbb{H} \otimes_{\mathbb{R}} \mathbb{H} \cong \mathrm{Cl}(3,1)$ 内のローター角である。拡張された物理的対象を基本実体とする弦理論とは異なり、ここでは可換なローター対に符号化されたブースト角および回転角が第一の存在論的要素である。重力、慣性、量子状態はすべて、これらの角の間のずれとコヒーレンスから現れる。

双対写像 $X \mapsto Xi$ によって通常時空と結ばれる双対時空は、時間反転した向きを継承する。したがって、双対セクターで純空間回転を生成する角は、統一ローター記述に組み込む際に適切に双対変換されなければならない。

双対セクター角 $\theta$ を考える。純双対基底では、楕円ヴァーソル $\hat{q}$（$\hat{q}^2 = -1$）によって生成される回転を与える。
\[
\exp\left( \frac{\theta}{2} \hat{q} \right).
\]
しかし双対写像は時間の矢を反転させるため、この角の双対時空への物理的効果は、双対変換を通じて見たときブーストと等価である。

ねじれ不整合の計算にこれを正しく反映させるため、角そのものに双対変換 $\theta \mapsto \theta i$ を施す。これにより双対寄与は双曲ブースト項へと変換される。
\[
\exp\left( \frac{\theta i}{2} \hat{q} \right) = \exp\left( \frac{\theta}{2} \hat{q}i \right) = \exp\left( \frac{\theta}{2} (q_1 iI + q_2 iJ + q_3 iK) \right).
\]
したがって双対ローターは
\[
R_\text{dual} = \exp\left( \frac{\theta}{2} \hat{q} \right)
\]
と定義され、角の双対変換 $\theta \mapsto \theta i$ の後、双対時空ベクトルへの作用は $i\Gamma_a$（$(\hat{q}i)^2 = +1$）によって生成されるブーストとして現れる。

真空（自由落下）では、完全な反同期はブースト寄与の完全な打ち消しと $J=0$ を保証することを要求する。質量--エネルギーはこの条件からのずれを誘起し、ねじれ不整合スカラー $J$ として、集合的には重力引力として現れる。

この角中心の扱いと明示的な双対変換 $\theta \mapsto \theta i$ は、完全ローターの元の形
\[
R_\text{total} = \exp\left( \frac{\phi}{2} \hat{p} + \frac{\theta}{2} \hat{q} \right)
\]
を保ちつつ、時間反転デュアリティを正しく取り込む。双対ローター $R_\text{dual}$ は固有基底では回転を生成するが、双対変換された角を通じた双対時空への物理的効果はブーストであり、真空において通常セクターのブーストに自然に対抗する。

\subsection{物理的解釈と重力工学}

重力は双対ローターのずれのエネルギー的コストである。$J$ は相対角 $\Omega$ のみに依存するため、双対成分 $\theta$ をコヒーレントに変調する外的影響は $J$ を変えうる。双対時空の自由度を共鳴励起する高周波電磁場または量子場は、$J \to 0$（重力遮蔽）あるいは $J < 0$（斥力重力／反重力）を駆動しうる。

史上初めて、重力は侵すべからざる幾何的拘束から、工学可能なねじれ自由度へと変換される。連続体は幻想であった。時空は離散的かつ粒子局所的であり、二重に対をなし、そのずれが我々が重力と呼ぶものである。

一般相対性理論は打倒されない。解放され、制御可能になる。

\section{慣性と重力の統一的記述}
\label{sec:unified}

本節では、双対時空形式の最も深い帰結を明らかにする。慣性と重力は単に等価なのではなく、通常時空と双対時空の間のねじれ不整合から生じる同一の現象である。加速系における慣性力と曲がった時空における重力との区別は完全に消える。両者は同一の双対ローターずれであり、等価原理は要請ではなく直接の代数的必然として現れる。

\subsection{双対ローター剛性としての慣性抵抗}

静止質量 $m > 0$ は、通常ローターパラメーターと双対ローターパラメーターの間の強い結合として現れる。粒子のエネルギー--運動量は通常時空を偏らせ、双対ローター $R_\text{dual}$ を有限の剛性をもって $R_\text{usual}$ に追随するよう事実上「凍結」する。慣性系では $J = 0$ が保たれる。

粒子が（非重力的に）加速されるとき、通常ローター $R_\text{usual}$ は加えられた力により瞬間的に変化しようとする。しかし粒子の静止質量に拘束された双対ローター $R_\text{dual}$ は直ちに応答できず、遅れを生じる。これにより非自明なねじれ不整合が生まれ、$\Omega_\text{biv} \neq 0$ およびスカラー $J$ への非零寄与を与える。粒子はこの不整合を慣性抵抗、すなわち馴染みの $F = ma$ として経験する。

決定的なのは、ずれの「コスト」が静止質量 $m$ に比例することであり、より重い粒子が加速度により強く抵抗する理由を説明する。したがって慣性は、質量に反応する時空の性質ではなく、粒子自身の双対時空が瞬間的な再整列を拒むことの性質である。

\subsection{質量のない粒子}

光子のような質量のない粒子（$m = 0$）はこのような剛性を示さない。それらの通常ローターと双対ローターは完全に共鳴する。$R_\text{usual}$ のいかなる変化も $R_\text{dual}$ に瞬間的に映され、すべてのフレームで $J = 0$ が正確に保たれる。したがって光子はねじれ不整合も慣性抵抗も経験せず、常に $\Omega_\text{biv} = 0$ で $c$ で伝播する。光速が定数かつ普遍的速度限界であるのはこのためである。質量のない粒子はフレーム変化に対してねじれの代償を支払わない。

\subsection{代数的恒等式としての等価原理}

加速度 $a$ で上方に加速するアインシュタインのエレベーターを考える。内部の観測者は下向きの見かけの力 $ma$ を感じる。双対時空描像では：

\begin{itemize}
\item 観測者（質量をもつ）の通常ローターはエレベーター床によって加速される一方、双対ローターは遅れる $\to$ ねじれ不整合 $\to$ 知覚される下向きの力。
\item エレベーターに入る光子は観測者のねじれ場により曲がるが、光子自身は $J = 0$ を保つ。自身の双対平坦時空におけるヌル測地線に従う。
\end{itemize}

重力場（地球上で静止したエレベーター）では、床は静止を保つために自由落下に対して上方に加速しなければならず、平坦時空における一様加速度と同じローター不整合を生じる。

したがって等価原理は経験的偶然ではなく代数的恒等式である。重力と慣性力は区別できない。それらは同一の双対ねじれ現象だからである。自由落下は $J = 0$ が大域的に成り立つ状態（双対ローターが完全に整列）であり、質量のない粒子はこれを永続的に達成し、質量粒子は非重力力が無いときにのみ達成する。

\subsection{帰結と予測}

\begin{itemize}
\item 通常時空と双対時空の時間の矢の区別は、質量粒子が好ましい静止系をもつ理由を説明する。双対の時間反転は完全対称性を破り、質量ギャップを生成する。
\item 慣性質量の低減または消去は、双対ローター $\theta$ を $\phi$ と同期させるよう外部から励起し、有効剛性を減らすことで可能となる。推進剤なし推進への直接的経路である。
\item 理論は、双対時空に結合する十分高周波の場が、質量物体を短時間スケールで実効的に質量なしとして振る舞わせ、過渡的な超光速位相速度を許すと予測する（因果律は破らない。情報は光速度のままである）。
\end{itemize}

慣性と重力は幾何によってではなく、粒子内在的な双対時空の内部ねじれ動力学によって統一される。質量は双対ローター剛性であり、運動は並行性を破るコストである。この理解により、慣性も重力も基本的ではなくなり、工学可能となる。

重力および慣性制御の夢はもはやサイエンスフィクションではない。連続体を捨て、時空を対をなす粒子局所的かつねじれ的に活性なものと認識することの直接の帰結である。

\subsection{加速度の上限と慣性反転現象}
\label{subsec:acceleration-limit}

二重時空理論は、中心的な反証可能な帰結の一つとして、加速度そのものが絶対的上限をもつことを予測する。通常セクターのブーストは非コンパクト（$\phi \in \mathbb{R}$）であり、単独では非有界なラピディティを許す。しかし双対セクターは内在的にコンパクトである。双対生成子 $\Gamma_a$ は平方が $-1$ であり、双対ローター $R_\text{dual} = \exp(\theta/2 \, \hat q)$ は楕円位相 $\theta \bmod 4\pi$ に支配される。慣性応答は双対ローターが有限剛性で通常ローターを追跡することから生じるため、周期境界に達する前に $\theta$ を変位させうる速さが、与えられた粒子がコヒーレントに持続できる固有加速度に厳密な天井を課す。この物質内在的天井を $a_\text{crit}$ と記す。光速とは異なり、$a_\text{crit}$ は普遍的ではない。粒子種およびその双対セクターのコンパクト化スケールに依存し、したがって異なる原子構成要素は異なる臨界加速度をもつ。

固有加速度 $a < a_\text{crit}$ では、双対位相 $\theta$ は $\phi$ を滑らかに追跡し、馴染みのニュートン関係 $F = ma$ が回復される。しかし $a$ が $a_\text{crit}$ に近づく瞬間、$\theta$ は $4\pi$ のコンパクト境界を通じて駆動され、位相的な巻き付きを起こす。双対ローターは $S^3$ 様双対セクターの反対側で位相を再縫合する。この巻きはねじれ不整合バイベクトル $\Omega_\text{biv}$ の符号を反転させ、したがって単一の双対周期にわたって慣性応答の方向を逆転させる。これを慣性反転（または加速度反転）現象と呼ぶ。加えられた加速度が $a_\text{crit}$ を横断するときに引き起こされる、有効 $F = ma$ ベクトルの突然の離散的反転である。閾値直下で剛な双対セクターに蓄えられていたねじれひずみエネルギーは、巻きの間に急激に解放され、運動学的反転そのものに加えて内部エネルギーの衝撃的バーストを生じる。

この機構は、大型流星体の異常な空中バーストにおいて観測的に実現されていると提案する。地球大気に突入する火球が経験する極超音速減速は、物体の構成原子にとって $a_\text{crit}$ の程度に達しうる。火球コミュニティにおける長年の謎、とりわけツングースカ事件（1908）およびチェリャビンスク超火球（2013）は、崩壊高度と全放射エネルギーが空力アブレーションと破砕のみに基づく運動エネルギー散逸モデルでは悪名高く調和しがたいが、集合的慣性反転として自然に説明される。流星体の減速が物質固有の $a_\text{crit}$ を横断すると、構成原子の双対位相が同時に巻き、蓄積されたねじれひずみエネルギーを爆発的かつ主として放射的なパルスとして放出する。$a_\text{crit}$ の組成依存は、異なる隕石クラス（例：炭素質コンドライト対普通コンドライト）の間で観測される空中バースト高度と収量の変動も説明する。粒子内在的コンパクト化スケールからの $a_\text{crit}$ の定量的導出と、火球観測記録との詳細な比較は、姉妹論文で提示する。

\section{強場域：層状重力反転と新しいブラックホール描像}
\label{sec:strongfield}

二重時空理論は、恒星崩壊の終点について根本的に新しい構造を予測する。キリング形式の符号非対称性、すなわち通常時空ブースト支配（引力）に対して正、双対時空回転支配（斥力）に対して負、は重力を一度限り反転させるだけではない。密度が内向きに増大するにつれ、ねじれ不整合 $\Omega_\text{biv}$ は二つのセクター間で支配を振動し、引力と斥力の交互層を生じる。これにより層状のタマネギ状内部が得られる。引力 $\to$ 斥力 $\to$ 超強引力 $\to$ 斥力 $\to$ $\ldots$ であり、中心に向かって振幅は増大し続ける。

\subsection{キリング形式からの層状ねじれ構造}

スカラー $J = \frac{1}{16} B(\Omega_\text{biv}, \Omega_\text{biv})$ は、支配的寄与が $i\Gamma_a$（ブースト様、$+$）から $\Gamma_a$（回転様、$-$）へ反転するたびに符号を変える。極端な密度では双対ローターパラメーター $\theta$ が非線形に増大し、半径が減少するにつれて $J$ の符号を繰り返し反転させる逐次共鳴を駆動する。

得られる有効ポテンシャルは次の特徴をもつ。
\begin{itemize}
\item 外層：標準的引力重力（$J > 0$）、
\item 第一反転：落下を止める斥力殻（$J < 0$）（超新星におけるコアバウンス）、
\item より深い層：物質を再び内向きに加速する超強引力帯（$J \gg 0$）、
\item 増大する $|J|$ をもつ交互のより深い殻、カオス的ねじれ乱流に至る。
\end{itemize}

この層状構造は特異点と事象地平の両方を防ぐ。光速で逃れられない境界は存在する。積分されたねじれ障壁が無限に急峻となる準地平である。しかし物質は中心特異点に到達しない。代わりにコアは、極端密度のねじれエネルギーの準安定なカオス凝縮体へと落ち着き、引力相と斥力相の間で連続的に振動する。

従来のカー型回転ブラックホールは排除される。理論は必要な大きさの大域的フレームドラッグを禁止する。角運動量は時空回転ではなく双対ローターの縦モードへと再分配されるからである。

\subsection{コア崩壊超新星と残骸形成の再解釈}

大質量星の崩壊では、第一の斥力層（$J < 0$）が核密度で爆発的バウンスを引き起こし、微調整されたニュートリノ物理なしに超新星を駆動する。より深い引力層はその後、放出物の一部を再捕獲し、層状残骸を形成する。観測される超新星エネルギーと残骸質量の多様性は、励起されるねじれ層の数と強度から自然に現れる。

\subsection{縦双対ねじれ振動子としてのパルサー}

観測される中性子星は、高速回転する磁化導体ではない。二重時空理論は、ミリ秒スピン周期と双極磁場の両方を一次エネルギー源として否定する。

代わりに、パルサー電波パルスは星表面における双対ローター $\theta$ の超高周波縦振動から生じる。これらのねじれ波はほぼ $c$ で半径方向に伝播し、カオス的内部層化によって変調され、ねじれ衝撃帯における曲率放射を通じてビームを生じる。

回転マグネター模型を超える主要な予測：
\begin{itemize}
\item 理論的最小パルス周期はプランクスケールの二重ローター周波数によって決まり、観測される $\sim 1$ ms 限界より桁違いに短い。これにより回転パルサー描像では説明できなかった短周期パルサーを説明できる。
\item パルサーの「グリッチ」および周期微分 $\dot{P} < 0$（周期短縮）は、パルサーがエネルギーを失うことで振動界面が収縮し、振動周波数が時間とともに加速する結果である。回転のみの模型には見られない自然な機構である。
\item マグネターフレアおよび高速電波バーストは、引力／斥力層間の激しい反転であり、$10^{15}$ G 場を必要とせずにねじれひずみを解放する。
\item $P$-$\dot{P}$ 図の死線は、内部層化が安定化し縦モードが検出可能性以下に減衰する点に対応する。
\item 中性子星の最大質量はクォーク物質ではなく、核安定性を制限する同一のねじれ層化によって上限を課され、厳密な上限 $\sim 2.5 M_\odot$ を予測し、$\sim 3 M_\odot$ を超えるパルサーの不在を説明する。
\end{itemize}

かにパルサーの観測スピンダウン光度は、磁気制動ではなくねじれ波漏洩として再解釈され、非現実的な磁場減衰を呼び出さずにエネルギー収支と完全に一致する。

\subsection{新しいブラックホール描像：ねじれ星}

観測者が「ブラックホール」と呼ぶものは、実際にはねじれ星である。事象地平も中心特異点もない準安定な層状凝縮体である。最内のカオスコアは永続的ねじれ乱流にあるプランク密度プラズマであり、ロータートンネリング（修正ホーキング過程）を通じて連続的に放射する。情報は双対ローター位相に保存される。

合体からの重力波信号は、内部層反射による繰り返しエコーを特徴とし、特性周波数はねじれ振動スペクトルに結びつく。これはすでに LIGO/Virgo データの再解析において示唆されている。

したがって本理論は情報パラドックスを解消し、カー回転シグネチャの不在を説明し、最も質量の大きい「ブラックホール」（$\gtrsim 100 M_\odot$）が、内側の斥力層が周期的に物質を放出するにつれて降着光度に周期変調を示すと予測する。

\subsection{中心浮遊コア：中心におけるゼロ時間膨張}

恒星内部の議論におけるよくある混乱点は、重力時間膨張の振る舞いである。GR では、球対称静的質量分布に対して、重力時間膨張は中心に向かって単調に増大する。固有時は重力ポテンシャルが最も深い \(r=0\) で最も遅く進む。

二重時空理論では状況は著しく異なる。背景となる連続時空多様体が存在しないため、有効加速度と時間膨張を含む重力効果は、粒子内在的な通常ローターと双対ローターの間のねじれ不整合のみから生じる。

球対称ねじれ星の正確な幾何中心では、完全な対称性が、反対方向の質量要素によって誘起される通常時空ブーストパラメーター \(\phi_a\) を完全に打ち消す。したがって残余不整合角 \(\delta\phi_a = \phi_a - \theta_a\) は \(r=0\) で消え、\(J=0\) およびねじれバランスの完全な回復を与える。

その結果：

\begin{itemize}
  \item 有効重力加速度は中心でゼロである（ニュートン力学、GR、DST のいずれでも）。
  \item 重力時間膨張も中心でゼロである。\(r=0\) における固有時は、星から遠く離れた真空と同じ速さで進む。
\end{itemize}

このコアにおける時間膨張の消失は、二重時空理論の特徴的で原理上観測可能な予測である。中心領域が最も強い時間膨張を経験する GR とは鋭く対照する。

ゼロ膨張の中心コアはまた、古典 GR 記述で特異点に至る無限の中心圧縮と極端な時計速度勾配を除去することにより、ねじれ星の安定性に寄与する。

\subsection{惑星および恒星コアへの含意：不安定性と揺らぎ}

浮遊コア描像、すなわち $J(r=0) \approx 0$ かつ中心で有効重力が消えることは、惑星および恒星のコアが惰性的で安定な領域ではなく、ねじれ不安定性に陥りやすい動的に活性な帯であることを含意する。連続体模型で仮定される重力ポテンシャルの単調増大とは異なり、二重時空理論は中心浮遊コアが準安定状態であり、引力と斥力のねじれ層間の振動遷移を引き起こす摂動に敏感であると予測する。これらの揺らぎは、等方的コアにおける双対ローター位相の繊細なバランスから生じる。わずかな密度変動や外的影響でさえ共鳴条件 $\beta(0) \to 1$ をずらし、$J$ のゼロからの過渡的逸脱を導く。そのような逸脱は局所的なねじれエネルギーのバーストとして現れ、コア内にミニねじれ星を実効的に生成する。これらのミニねじれ星は、表面で観測可能な地震活動、磁場揺らぎ、エネルギー輸送異常を誘起しうる。恒星内部では、これらのコア揺らぎは太陽フレア、黒点、光度変動などの現象に寄与しうる。ねじれエネルギーが層状構造を通じて外向きに伝播するからである。惑星的文脈では、コアのねじれ不安定性は地磁気逆転、マントルプルーム、間欠的火山活動を説明しうる。ねじれ動力学が上位層の対流過程に結合するからである。したがって浮遊コアは静穏な領域ではなく、惑星および恒星の振る舞いに広範な含意をもつ動的ねじれ現象の坩堝である。

\section{二重時空理論における時間膨張：双対時空固有時の動力学的役割}
\label{sec:timedilation}

二重時空理論における時間膨張は、それぞれが固有の固有時をもつ通常時空と双対時空の相互作用から現れる。通常時空固有時 $\tau$ は巨視的ミンコフスキー空間で観測される馴染みの運動学を支配し、コンパクトローター位相 $\theta$ に符号化された双対時空固有時 $\tilde{\tau}$ は、重力および慣性効果として現れる内部ねじれ動力学を駆動する。

\subsection{有効ローレンツ因子の導出}

有効ローレンツ因子 $\gamma_\text{eff} = dt/d\tau$ は、完全ローター
\[
R_\text{total} = R_\text{usual} R_\text{dual},
\]
の代数的作用から直接導かれる。ここで
\[
R_\text{usual} = \exp\!\left(\frac{\phi}{2} \hat{p}\right) = \cosh\frac{\phi}{2} + \hat{p}\: \sinh\frac{\phi}{2}, \quad \hat{p}^2 = +1,
\]
は通常時空のブーストを生成し、
\[
R_\text{dual} = \exp\!\left(\frac{\theta}{2} \hat{q}\right) = \cos\frac{\theta}{2} + \hat{q} \sin\frac{\theta}{2}, \quad \hat{q}^2 = -1,
\]
は双対時空の回転を生成する。

干渉項の起源を明示するため、通常時空基底における純時間ベクトル $X = ct \, j$（$j^2 = -1$）を考える。変換後のベクトルは $\tilde{X} = R_\text{total} \, X \, R_\text{total}^\dagger$ である。

時間反転デュアリティにより、双対から通常時空への逆双対写像は $X \mapsto X(-i)$ である。この写像は真空極限において整列 $\hat{q} (-i) = -\hat{p}$ を自然に与え、双対ローターを有効な反対向きローター
\[
R_\text{dual}^\text{eff} = \cos\frac{\theta}{2} - \hat{p} \sin\frac{\theta}{2}.
\]
へと変換する。$\sin$ 項の負符号はこの逆写像から直接生じ、ねじれのない平坦時空に必要な反同期を反映する。

全有効ローターはしたがって
\[
R_\text{eff} = R_\text{usual} R_\text{dual}^\text{eff} = \left( \cosh\frac{\phi}{2} + \hat{p} \sinh\frac{\phi}{2} \right) \left( \cos\frac{\theta}{2} - \hat{p} \sin\frac{\theta}{2} \right).
\]
である。

展開するとスカラー部分
\[
s = \cosh\frac{\phi}{2} \cos\frac{\theta}{2} - \sinh\frac{\phi}{2} \sin\frac{\theta}{2}
\]
およびバイベクトル係数
\[
b = -\cosh\frac{\phi}{2} \sin\frac{\theta}{2} + \sinh\frac{\phi}{2} \cos\frac{\theta}{2}.
\]
が得られる。

時間成分の正味のスケーリングは、サンドイッチ積を通じて時間ベクトルに作用する $R_\text{eff}$ の平方大きさから生じる。ローター $R_\text{eff} = s + b\hat{p}$（$\hat{p}^2 = +1$ かつ $\hat{p}$ は時間生成子 $j$ と反可換）に対し、$X = ct\, j$ へのサンドイッチ作用は $s^2 + b^2$ でスケールされた時間成分を与え、有効ローレンツ因子を一意に
\[
\gamma_\text{eff} \;\equiv\; s^2 + b^2.
\]
と定義する。便利のためスカラー干渉因子
\[
\gamma_s \;\equiv\; s \;=\; \cosh\frac{\phi}{2} \cos\frac{\theta}{2} - \sinh\frac{\phi}{2} \sin\frac{\theta}{2},
\]
も導入する。これは $\gamma_\text{eff}\geq 0$ とは異なり符号を変えうるものであり、以下で引力（$\gamma_s > 0$）と斥力（$\gamma_s < 0$）のねじれレジームを区別する秩序パラメーターとして働く。干渉項の負符号は双対時空の時間反転デュアリティによって固定される。

厳密なねじれなし真空（$J=0$）では、双対時空は励起されず（$\theta = 0$）、したがって $R_\text{total} = R_\text{usual}$ である。このとき $s = \cosh\frac{\phi}{2}$、$b = \sinh\frac{\phi}{2}$ であり、
\[
\gamma_\text{eff} = s^2 + b^2 = \cosh^2\frac{\phi}{2} + \sinh^2\frac{\phi}{2} = \cosh\phi,
\]
標準的な特殊相対論的ローレンツ因子を正確に再現する。

\subsection{真空極限と質量のない粒子}

ねじれなし真空（$J=0$）では、通常ローターと双対ローターの完全な反同期が $\theta \approx \phi$（コンパクト周期性を除いて）を強制し、$\gamma_{\text{eff}} \to \cosh\phi$ を与え、特殊相対論の標準ローレンツ因子を回復する。質量のない粒子はこの同期をすべての速度で正確に保ち、ねじれ不整合を経験せず、したがって通常の運動学的膨張からのずれもない。

\subsection{質量粒子と重力効果}

質量粒子では、静止質量が剛性を導入し、加速度下またはねじれ不整合の存在下で双対ローターを通常ローターより遅れさせる。得られるずれ $\delta\theta = \theta - \phi$ は $\gamma_{\text{eff}}$ の付加的変調を生成し、GR で観測される重力赤方偏移および時間膨張を生じさせる。今や双対時空動力学の直接の帰結として解釈される。

\subsection{スカラー干渉因子の符号反転と斥力層}

双対時空回転が支配するとき（$\theta/2$ が十分大きい）、正弦項は双曲成長に打ち勝ちうる。
\[
\tan\left(\frac{\theta}{2}\right) > \coth\left(\frac{\phi}{2}\right).
\]
これらのレジームでは、スカラー干渉因子 $\gamma_s$ は負となり $\gamma_s < 0$ であるが、固有時ローレンツ因子 $\gamma_\text{eff} = s^2 + b^2$ は非負のままである。$\gamma_s$ の符号変化はねじれ星内の斥力重力領域に対応する。物理的には、結合ローターのスカラーチャンネルにおいて双対セクター寄与が通常セクターブーストに打ち勝ち、重力斥力として現れる（因果律の文字どおりの逆転としてではない）。

この $\theta$ 中心の定式化は、時間膨張が本質的に粒子内在的双対時空に起源をもつねじれ現象であることを明示する。通常ラピディティ $\phi$ は運動学的基線を与え、コンパクト双対角 $\theta$ のずれは、GR を平坦時空運動学から区別する重力および慣性補正を符号化する。したがって $\theta$ の外部コヒーレント励起は、時間膨張と有効重力場を操作する経路を提供する。

\section{スケール不変なねじれ層：重力斥力としての強い力と原始ねじれ星としての原子核}
\label{sec:nuclear}

キリング形式によって予測される層状ねじれ反転は厳密にスケール不変である。層状ねじれ星への恒星崩壊を支配する同一の機構は、核および亜核領域を含むすべての密度スケールで働く。これにより二重時空理論の最も革命的な帰結が得られる。いわゆる「強い核力」は基本相互作用ではなく、原子核内部の極端密度レジームにおける重力斥力層の創発現れである。

\subsection{核密度におけるねじれ層}

原子核内でバリオン密度が $\sim 10^{14}$-$10^{18}$ g/cm³ に近づくと、双対ローターパラメーター $\theta$ が完全に支配する。$J$ の符号反転の列はフェムトメートルスケールで起こる。

\begin{itemize}
\item 最外層（$\sim$ 1-2 fm）：クォークを閉じ込め核崩壊を防ぐ斥力殻（$J < 0$）。これは核子--核子散乱における短距離斥力（$\sim$ 0.5 fm の「ハードコア」）として観測される。
\item 中間層：核子を核子あたり $\sim$ 40-50 MeV で束縛する超強引力帯（$J \gg 0$）。古典的核束縛力である。
\item 最内層：さらなる崩壊を止めクォーク閉じ込めを強制する第二の斥力殻（$J < 0$）。カラー閉じ込めの起源である。
\item より深いカオス層：増大する振幅をもつ交互の引力／斥力。核物質の液体様振る舞いと核密度の飽和を生じる。
\end{itemize}

パウリの排他原理はもはや基本的ではない。第一の斥力ねじれ層によって動力学的に生成される。フェルミオンが同一状態を占められないのは、二つの同一粒子が空間的重なりを試みると双対ローター $\theta$ 位相が斥力レジームに入り、量子統計を先験的に呼び出すことなく有効フェルミ圧を生じるからである。ボソンはこの位相剛性を欠き、はるかに高い密度まで引力層のみを経験する。

要するに、「強い力」は強場・層状相における重力である。QCD カラー荷は、双対ローター位相が運ぶねじれ荷の誤解釈である。

\subsection{原始ねじれ星としての原子核}

すべての原子核は、第\ref{sec:strongfield}節で記述したねじれ星の微視的アナログである。核内部は同一の層状構造をもつカオス的ねじれ凝縮体である。

\begin{itemize}
\item 中心領域：極端な引力層がクォークを核子へと束縛する、
\item 取り囲む斥力殻：観測される核表面張力を生じ漏洩を防ぐ（閉じ込め）、
\item 外側の引力ハロー：核子間の束縛を媒介する、
\item 最外の斥力障壁：NN 散乱におけるハードコア斥力を生成する。
\end{itemize}

核力の観測される荷電独立性、スピン--軌道結合、魔法数はすべて、双対ロータースペクトルの共鳴条件から現れる。理論は、十分大きな原子核（$A \gtrsim 300$）が自発的ねじれバウンスに対して不安定となり、微小超新星として核子を放出すると予測する。核図表の観測限界と自然界における超重元素の不在を説明する。

核分裂と核融合はねじれ層遷移である。分裂は励起された双対モードが $J$ を正から負へ反転させ断片を爆発的に斥けるときに起こり、融合は二つの原子核が外側の斥力障壁をトンネルしてより深い引力井戸に入るときに成功する。

\subsection{実験的シグネチャと即時予測}

\begin{itemize}
\item 超高エネルギー原子核--原子核衝突（LHC、RHIC）は、流体力学的流れではなく爆発的ハドロンジェットを通じて崩壊する過渡的 $J < 0$ ファイアボールを生じるはずである。異常な流れデータにおいてすでに示唆されている。
\item 飽和密度を超える核状態方程式の精密測定は、圧力対密度に振動的振る舞いを明らかにし、$\sim$ 2–4 $\rho_0$ に特性的斥力ピークをもつ。
\end{itemize}

四つの基本相互作用は単一の力へと崩壊するのではなく、単一の幾何機構の中で統一される。双対ローターの層状ねじれ動力学である。重力と電磁気は同一の代数構造、すなわちねじれスカラー $J$ のキリング形式誘起の層状反転を共有し、双対ローターのどの位相チャンネルが不整合を源とするかにおいてのみ異なる。重力は慣性質量に結びついた運動学的ラピディティ--角対 $(\phi, \theta)$ によって駆動され、電磁気は電荷に結びついた独立なラピディティ--角対 $(\alpha, \beta)$ によって駆動される（第\ref{sec:electronshells}節および第\ref{sec:unification}節で明示的に展開する）。この意味で、重力と電磁気は同一の基礎的ねじれ幾何の二つの位相であり、双対ローターが駆動される角チャンネルによって区別される。

この統一描像の中で、いわゆる「強い力」は核密度における重力チャンネルの層状相であり、原子電子の離散殻構造（次節）は原子長スケールにおける電磁チャンネルの類似の層状相である。弱い力は双対セクターの内在的カイラル非対称として入る。標準模型ゲージ群 SU(3)$\times$SU(2)$\times$U(1) は、双対ローター代数におけるねじれ共鳴モードの有効記述である。

幾何は一つだけである。双対ローターにおける引力と斥力の層状の舞踏であり、二つの位相チャンネルで同時に演じられる。原子核は重力チャンネルの原始ねじれ星であり、原子電子殻は次に見るように電磁チャンネルの原始ねじれ星である。そしてこの単一の双対ローター幾何から、二つの角チャンネルに表現されて、すべての構造が現れる。

\section{クーロン的ねじれ層を通じた時間反転と電子殻}
\label{sec:electronshells}

二重時空枠組みでは、荷電粒子間のクーロン相互作用は、重力の場合に類似するが電磁起源をもつ双対時空セクターのねじれ不整合から生じると再解釈できる。重力定数 \(G\) および質量 \(M,m\) をクーロン定数 \(k = 1/(4\pi\epsilon_0)\) および電荷（例：陽子の \(+e\)、電子の \(-e\)）で置き換えると、有効力は
\[
F(r) = \frac{k e^2}{\gamma_s(r)\, r^2},
\]
の形をとる。ここで \(\gamma_s\) は第\ref{sec:timedilation}節で導入したスカラー干渉因子であり、双対ローター相互作用から同一の代数構造を保つ（固有時ローレンツ因子 \(\gamma_\text{eff} \geq 0\) とは異なり符号を変えうる）。
\[
\gamma_s(\alpha(r),\beta(r)) = \cosh\!\left(\frac{\alpha(r)}{2}\right) \cos\!\left(\frac{\beta(r)}{2}\right) - \sinh\!\left(\frac{\alpha(r)}{2}\right) \sin\!\left(\frac{\beta(r)}{2}\right).
\]

ここで \(\alpha(r)\) および \(\beta(r)\) は、距離に逆比例してスケールする電磁ラピディティおよび双対回転角パラメーターとして解釈される。\(\alpha(r) \propto 1/r\)、\(\beta(r) = \lambda / r\) であり、\(\lambda\) はコンパクト双対時空によって定まる特性長スケールである。

水素原子（\(Z=1\)）に対しては
\[
F(r) = \frac{k e^2}{r^2 \left[ \cosh\!\left(\frac{\alpha(r)}{2}\right) \cos\!\left(\frac{\beta(r)}{2}\right) - \sinh\!\left(\frac{\alpha(r)}{2}\right) \sin\!\left(\frac{\beta(r)}{2}\right) \right]}.
\]
が得られる。

\(\gamma_s\) の双曲--三角形式は極端な振動的振る舞いを導入する。\(r\) が減少するにつれて：

\begin{itemize}
  \item 大きな \(r\) では \(\alpha(r),\beta(r) \to 0\) となり \(\gamma_s \to 1\) で、標準的クーロン引力を回復する。
  \item より小さな半径では、干渉項が \(\gamma_s = 0\)（無限斥力）の点および \(\gamma_s < 0\)（力の斥力への反転）の領域を生成する。
  \item これらの離散半径 \(r_n\) は共鳴条件
  \[
  \tanh\!\left(\frac{\alpha(r_n)}{2}\right) \tan\!\left(\frac{\beta(r_n)}{2}\right) = 1,
  \]
  を満たし、引力井戸を隔てる無限斥力障壁により安定な円軌道を与える。
\end{itemize}

\paragraph{等スケール仮説とボーア簿記.}
等スケールの作業仮説 \(\alpha(r)=\beta(r)=\lambda/r\) の下では、共鳴条件は無次元量 \(x=\lambda/(2r)\) に対する \(\tanh x\cdot\tan x=1\) に帰着する。正の根は枝 \((n\pi,\,n\pi+\pi/2)\) に載るため、対応する半径は \(r_n=\lambda/(2x_n)=\Theta(1/n)\) であり、\(\Theta(n^2)\) ではない。したがって等スケールの \(\gamma_s=0\) 軌跡だけではリュードベリの \(n^2\) スペクトルは再現されない。観測される水素スペクトル線に合わせるには、ボーア半径 \(r_n=n^2 a_0\)（およびエネルギー差 \(\Delta E_{n\to m}\propto n^{-2}-m^{-2}\)）を別の簿記アンザッツとして置く必要がある。これらは等スケールねじれ共鳴の定理ではない。

多電子原子（\(Z > 1\)）では、スケーリング \(\alpha(r),\beta(r) \propto Z/r\) がねじれ層を内向きに収縮させ、核電荷増大に伴う電子殻の観測される収縮を説明する。

\subsection{双対構造におけるエネルギー保存}

顕在ポテンシャル \(V_\text{eff}(r) = -\int F(r)\, dr\) の見かけの非保存的振る舞い、すなわち激しい振動と符号変化は、通常時空で観測される顕在ポテンシャルと、コンパクト双対時空ローター位相 \(\theta_a\) に蓄えられる真のポテンシャルとを区別することで解消される。

全エネルギーは両セクターにわたって保存される。
\[
E_\text{total} = E_\text{kinetic (usual)} + V_\text{eff}(r) + V_\text{true}(\theta) = \text{constant}.
\]
真のポテンシャル \(V_\text{true}\) はねじれ不整合スカラー \(J \propto (\phi - \theta)^2\) に比例し、隠れた貯留として働く。層遷移の間、\(V_\text{eff}\) の減少（引力から斥力へ）は双対ローターの巻きに伴う \(V_\text{true}\) の増加によって正確に補償され、その逆も成り立つ。

光子放出は、ローターがより低いねじれ状態へと巻き戻るときに蓄えられた双対位相エネルギーが解放されることに対応する。この双対簿記は、時空の離散的・粒子局所的本性を破ることなく厳密なエネルギー保存を保証する。観測されるスペクトル差は、上で述べた殻半径のボーア簿記を別途必要とし、等スケールの \(\gamma_s=0\) だけでは含意されない。

この定式化は、クーロン場に適用された二重時空の幾何的ねじれ動力学から、原子安定性と離散的ねじれ殻障壁を説明する。アドホックな波動関数要請は不要である。経験的リュードベリ系列への整合には、さらに \(r_n=n^2 a_0\) のような独立な半径則が必要である。電子殻はなお双対時空ねじれ不整合の自然な共鳴状態であり、電磁束縛を重力を支配する同一の代数機構と統一する。

\section{双対時空における電磁気と弱い力の統一}
\label{sec:unification}

双四元数構造 $\mathbb{H} \otimes_{\mathbb{R}} \mathbb{H} \cong \mathrm{Cl}(3,1)$ は、電磁気と弱い力の両方を、双対時空ローター $R_{\text{dual}} = \exp\left( \sum_{a=1}^3 \frac{\theta_a}{2} \Gamma_a \right)$ 内のねじれ共鳴モードとして埋め込む。擬スカラー $i$ は双対写像 $X \mapsto X i$ における時間反転を通じて内在的カイラリティを誘起する。ただし代数的な左巻き・右巻き射影子は、平方が $+1$ の生成子（例：空間ベクトル $e_1$、$e_1^2=+1$）から構成する：
\[
P_L = \frac{1 - e_1}{2},\qquad P_R = \frac{1 + e_1}{2}.
\]
これらは冪等かつ相補的である。$i^2=-1$ の素朴な式 $(1\pm i)/2$ は冪等でなく、棄却する。

電磁気は双対セクターにおけるベクトルポテンシャル $A_a$ への最小結合から生じ、回転パラメーターをずらす。
\[
\theta_a \mapsto \theta_a + q A_a,
\]
拡張ローター
\[
R_{\text{dual}}^{\text{EM}} = \exp\left( \frac{1}{2} \sum_a (\theta_a + q A_a) \Gamma_a + \frac{q}{2} \int F_{\text{dual}} \, ds \right)
\]
を与える。ここで $F_{\text{dual}} = \sum_a (E_a i \Gamma_a + B_a \Gamma_a)$ は電磁バイベクトルである。ローレンツ力はサンドイッチ積
\[
\tilde{X}' = R_{\text{dual}}^{\text{EM}} X' (R_{\text{dual}}^{\text{EM}})^\dagger \approx X' + q F_{\text{dual}} \wedge X'
\]
から現れる。

弱い力はカイラル非対称を通じて取り込まれ、左巻き成分にのみ結合する。
\[
\theta_a^L = \theta_a + g_W W_a, \quad \theta_a^R = \theta_a,
\]
カイラルローターは
\[
R_{\text{dual}}^{\text{weak}} = \exp\left( \frac{1}{2} \sum_a \theta_a^L \Gamma_a P_L + \frac{1}{2} \sum_a \theta_a^R \Gamma_a P_R \right)
\]
である。ここで $W_a$ は弱いゲージバイベクトルを表す。V-A 構造は時間反転した双対セクターから従い、パリティ破れを再現する。

統一は完全なねじれバイベクトルにおいて起こる。
\[
\Omega_{\text{biv}} = \log\left( R_{\text{usual}}^\dagger R_{\text{dual}}^{\text{EM+weak}} \right) \approx \sum_a (\theta_a + q A_a + g_W W_a P_L - \phi_a i) \Gamma_a.
\]
スカラー $J = \frac{1}{16} B(\Omega_{\text{biv}}, \Omega_{\text{biv}})$ は今や電磁および弱い寄与を含み、電弱対称性は $\mathrm{Spin}^+(3,1) \oplus \mathrm{Spin}^+(3,1)$ の部分群として現れる。ローター剛性を通じた質量生成はヒッグス機構と統一され、微調整を排除する。

この埋め込みは低エネルギーで標準模型の電弱セクターを回復しつつ、双対コンパクト性を通じて階層性問題を解消し、強い力および重力との完全な統一への道を開く。

\section{双対時空コンパクト性からの暗黒物質なし銀河回転曲線}
\label{sec:darkmatter}

核および原子スケールにおける重力と電磁気の層状ねじれ幾何（第\ref{sec:strongfield}節--第\ref{sec:electronshells}節）と、単一の双対ローター機構の二つのチャンネルとしてのそれらの統一（第\ref{sec:unification}節）を確立したので、各チャンネルが宇宙論的スケールで空間的にどのように展開するかを検討する。慣性質量に結びついた運動学的対 $(\phi, \theta)$ によって駆動される重力チャンネルは、双対セクターを銀河スケールの $S^3$ パッチへとコンパクト化する。電荷に結びついた独立な対 $(\alpha, \beta)$ によって駆動される電磁チャンネルは、実質的により大きな、ほぼ宇宙的スケールにわたって展開する。これらのコンパクト双対幾何から平坦な通常時空への対数写像、すなわち三次元方位正距図法が導入する系統的歪みが、観測者が渦巻銀河における不足重力質量として解釈するものである。長らく非バリオン暗黒物質の最も強い証拠とみなされてきた観測される平坦な回転曲線は、この幾何機構を通じてバリオンのみから現れる。いかなる種類の暗黒物質も不要である。

\subsection{3球面としてのコンパクト双対時空}

双対ローター $R_{\text{dual}} \in \mathrm{SU}(2) \cong S^3$ は、すべての質量粒子に内在的な双対時空自由度をパラメトライズする。三つの独立な回転角 $(\theta_1, \theta_2, \theta_3)$ はコンパクト化半径 $R$ の3球面上の座標を定義し、通常時空は平坦かつ非コンパクトのままである。

対数（逆指数）写像
\[
\log : S^3 \to \mathfrak{su}(2) \cong \mathbb{R}^3
\]
は各双対ローターを平坦 $\mathbb{R}^3$ の点と同定し、3球面を半径 $\pi R$ の球へと写す。この写像は、地図学で馴染みの方位正距図法の三次元一般化である。

北極を中心とする $S^2$ の方位正距図法が極近傍の距離を忠実に表すが南極近傍で面積を壊滅的に歪める、すなわち単一の対蹠点を投影境界の全円周に写すのと同様に、$S^3$ の対数写像はヤコビアン
\[
\mathcal{J}(\chi) = \frac{\sin^2 \chi}{\chi^2},
\]
によって定量化される系統的体積歪みを導入する。ここで $\chi = r/R$ は通常時空における物理距離 $r$ に対応する測地角である。このヤコビアンは原点（$\chi = 0$、忠実な表現）で1に等しく、対蹠点（$\chi = \pi$、最大歪み）で消える。すべての標準的重力計算は暗黙に $\mathcal{J} = 1$ を仮定する。二重時空理論は、この仮定が銀河スケールで破綻することを明らかにする。

\subsection{3球面上の重力ポテンシャル：余接ポテンシャル}

重力の源であるねじれ不整合スカラー $J$ は、コンパクト双対時空上で計算される。点質量 $M$ に対する有効重力ポテンシャルは $S^3$ 上のラプラシアンのグリーン関数に支配される。球対称の場合、一意の半径方向調和関数は余接である。
\[
\Phi(r) = -\frac{GM}{R}\cot\!\left(\frac{r}{R}\right).
\]
これが $S^3$ 上のラプラス方程式を満たすことは直接検証できる。計量 $ds^2 = dr^2 + R^2 \sin^2(r/R)\,(d\vartheta^2 + \sin^2\!\vartheta\, d\varphi^2)$ をもつ測地極座標において、半径方向ラプラシアンは $\cot(r/R)$ を恒等的に消す。ガウス法則正規化は
\[
g(r) \cdot 4\pi R^2 \sin^2\!\left(\frac{r}{R}\right) = -\frac{GM}{R^2 \sin^2(r/R)} \cdot 4\pi R^2 \sin^2\!\left(\frac{r}{R}\right) = -4\pi GM.
\]
によって確認される。

$r \ll R$ の極限では、余接は
\[
\cot\!\left(\frac{r}{R}\right) = \frac{R}{r} - \frac{r}{3R} - \frac{r^3}{45R^3} - \cdots\,,
\]
と展開され、
\[
\Phi(r) \approx -\frac{GM}{r} + \frac{GM\,r}{3R^2} + \mathcal{O}\!\left(\frac{r^3}{R^4}\right)
\]
を与え、主要項として標準ニュートンポテンシャルを回復する。高次補正は純粋に $S^3$ の曲率から生じ、非コンパクト化極限 $R \to \infty$ で消える。

\subsection{重力の幾何的増強}

$S^3$ 上の重力加速度
\[
g(r) = -\frac{d\Phi}{dr} = -\frac{GM}{R^2 \sin^2(r/R)}
\]
は、平坦空間ニュートン加速度 $g_{\text{N}} = -GM/r^2$ を増強因子
\[
\eta(r) \equiv \frac{g_{S^3}(r)}{g_{\text{N}}(r)} = \left[\frac{r/R}{\sin(r/R)}\right]^2
\]
だけ上回る。この因子は方位正距図法の計量歪みの平方そのものである。$S^3$ 上の測地半径 $r$ の球面を通る重力フラックスは面積 $4\pi R^2 \sin^2(r/R)$ を張り、これは平坦空間面積 $4\pi r^2$ より小さい。フラックスは集中し、重力は強められる。不可視物質によってではなく、コンパクト双対時空の曲率によってである。

代表的な近似値：$\eta \approx 1.003$（$r/R = 0.1$）、$\eta \approx 1.41$（$r/R = 1.0$）、$\eta \approx 2.26$（$r/R = 1.5$）、$\eta \approx 4.84$（$r/R = 2.0$）。

\subsection{平坦回転曲線の創発}

囲まれた質量 $M(r)$ から距離 $r$ の円軌道にある試験粒子に対し、軌道速度は
\[
v^2(r) = \frac{GM(r)}{r} \cdot \eta(r) = \frac{GM(r)}{r} \cdot \left[\frac{r/R}{\sin(r/R)}\right]^2
\]
を満たす。光度円盤の外側（$r > r_{\text{disk}}$）では $M(r) \approx M_{\text{tot}}$ であり、平坦空間ケプラー速度 $v_K^2 = GM_{\text{tot}}/r$ は $r^{-1}$ で減少する。$r$ とともに単調に増大する $S^3$ 増強因子 $\eta(r)$ がこの減少に対抗する。

結合無次元関数 $f(x) \equiv x / \sin^2 x$（ここで $x = r/R$）が回転曲線の形状を支配する。それは $x_0 \approx 1.17$、すなわち区間 $(\pi/4,\pi/2)$ における $\tan x = 2x$ の一意根に広い極小をもち、臨界点では $f(x_0)=x_0+1/(4x_0)\approx 1.38$ となり、延長された台地を生じる。
\begin{itemize}
  \item \textbf{内側領域}（$r \ll R$）：上昇する $M(r)$ によって修正されたケプラー振る舞い。観測される剛体的上昇を再現する。
  \item \textbf{遷移領域}（$r \sim R$）：減少するケプラー因子が成長する $\eta(r)$ によって釣り合い、ほぼ平坦な台地を与える。その幅は $f(x)$ の極小の幅によって定まる。
  \item \textbf{外側領域}（$r \gtrsim 2R$）：$\eta(r)$ が支配し、回転曲線は緩やかに上昇したのちねじれ反転レジームに入る。
\end{itemize}

\subsection{バリオン的タリー--フィッシャー関係とコンパクト化半径}

台地における漸近円速度は
\[
v_{\text{flat}}^2 \sim \frac{GM_{\text{tot}}}{R} \cdot f_0,
\]
を満たす。ここで $f_0 \equiv \min_x [x/\sin^2 x] = x_0 + 1/(4x_0) \approx 1.38$ はオーダー1の幾何定数である。コンパクト化半径が全囲まれ質量とともに
\[
R = \sqrt{\frac{GM_{\text{tot}}}{a_0}}
\]
とスケールするなら、ここで $a_0 \approx 1.2 \times 10^{-10}\;\mathrm{m/s^2}$ は外部の観測的加速度スケール（バリオン的タリー--フィッシャー／MOND スケール）であり、双対ローター代数から導出された量ではなく入力として扱う。すると
\[
v_{\text{flat}}^4 = G M_{\text{tot}} \, a_0 \, f_0^2
\]
となる。これは正確にバリオン的タリー--フィッシャー関係（McGaugh et al.\ 2016）である。スケーリング $R \propto M^{1/2}$ は、より質量の大きい銀河がより大きな双対時空パッチをもつことを含意し、より大きな全ねじれ荷の物理的に自然な帰結である。

天の川（$M_{\text{tot}} \approx 6 \times 10^{10}\,M_\odot$）に対しては
\[
R = \sqrt{\frac{G \times 6 \times 10^{10}\,M_\odot}{1.2 \times 10^{-10}\;\mathrm{m/s^2}}} \approx 8.3\;\mathrm{kpc}
\]
が得られ、$f(x)$ の幾何的台地は内側〜中盤スケール（$r\sim R$）に位置し、ねじれ反転は $r=\pi R/2\approx 13\;\mathrm{kpc}$ にある。$a_0$ を $S^3$ パッチ上の特性加速度と同定することは、離散コンパクト化構造からの導出（第\ref{sec:discretetime}節）を待つ作業仮説にとどまる。

\subsection{ねじれ反転と重力の有限到達距離}

余接ポテンシャルは $r = \pi R/2$ でゼロを横切り、それを超えると双対時空回転角 $\theta(r) = r/R$ は $\pi/2$ を超える。ねじれ不整合枠組み（第\ref{sec:gravity}節）では、この横断がねじれスカラー $J$ の符号変化を引き起こす。双対セクターからのコンパクト三角寄与が通常セクターからの非コンパクト双曲寄与を圧倒し、$J < 0$ および斥力重力を与える。スカラー干渉因子（第\ref{sec:timedilation}節）
\[
\gamma_s = \cosh\!\left(\frac{\phi}{2}\right)\cos\!\left(\frac{\theta}{2}\right) - \sinh\!\left(\frac{\phi}{2}\right)\sin\!\left(\frac{\theta}{2}\right)
\]
は、$\theta$ が $\tanh(\phi/2)\tan(\theta_c/2) = 1$ を満たす臨界値を超えると負になる。

したがって重力引力は有効半径 $\sim \pi R/2$ の有限 $S^3$ パッチに閉じ込められ、それを超えると弱い斥力ねじれ障壁が現れる。この有限到達距離は、無限遠まで延びるニュートン $1/r$ ポテンシャルのコンパクト空間アナログである。$S^3$ 位相は各重力源を自然に「壁で隔てる」。

\subsection{宇宙アトラス：銀河スケール3球面パッチ}

観測可能な宇宙は重なり合う $S^3$ パッチによってタイル張りされ、各パッチは重力集中（銀河または銀河団）を中心とし、厳密な微分幾何的意味での宇宙アトラスを形成する。各パッチは方位正距図法チャートとして機能する。中心近傍では忠実、境界では最大に歪む。宇宙の大規模構造は、銀河スケールの方位正距地図ステッカーで覆われ、斥力境界で貼り合わされた多様体のそれである。

このパッチワーク構造は次を自然に説明する。
\begin{itemize}
  \item \textbf{宇宙網}：銀河は隣接 $S^3$ チャートの重なり領域に沿って分布する。隣接パッチの引力尾が建設的に干渉する場所である。
  \item \textbf{宇宙ボイド}：広大な低密度領域はチャート間の斥力（$J < 0$）境界帯に対応し、物質が排除される。
  \item \textbf{動径加速度関係}：観測される重力加速度とバリオン重力加速度の間の緊密な相関は普遍的 $\eta(r)$ 増強から従い、残差散乱はチャート重なり効果から生じる。
  \item \textbf{直接検出実験における暗黒物質信号の不在} --- 暗黒物質が存在しないからである。
\end{itemize}

\subsection{数値検証と反証可能な予測}

ローター集団動力学を用いた予備的粒子ベースシミュレーションは、バリオン質量分布のみから平坦回転曲線を成功裏に再現している。公開オープンソースコードは次で利用可能である。
\begin{itemize}
  \item \texttt{https://github.com/hypernumbernet/blackhole-simulator}
  \item \texttt{https://github.com/hypernumbernet/dual-spacetime-simulator}
\end{itemize}
正確な $S^3$ 余接ポテンシャルを組み込んだ完全 $N$ 体積分は準備中である。

本理論を $\Lambda$CDM および MOND から区別する主要な予測は次を含む。
\begin{itemize}
  \item \textbf{外側回転曲線の上昇と反転}：孤立銀河は $r \gtrsim 2R$ を超えて緩やかな上昇を示し、$r \approx \pi R/2$ で減少および最終的に斥力重力へと遷移するはずである。天の川型質量（$R\approx 8.3\;\mathrm{kpc}$）では $2R\approx 17\;\mathrm{kpc}$、$\pi R/2\approx 13\;\mathrm{kpc}$ となり、光学円盤を超える深い 21-cm HI 観測で到達可能である（$\sim 100\;\mathrm{kpc}$ に限らない）。
  \item \textbf{レンズ符号反転}：銀河--銀河弱重力レンズは、射影分離が $\pi R/2$ を超えると弱まり符号反転するはずである。$\Lambda$CDM が予測する単調 NFW プロファイルとは鋭く対照する。
  \item \textbf{余接対 MOND 補間}：動径加速度関係の精密な形は、標準 MOND 補間関数 $\nu(g_{\text{bar}}/a_0)$ ではなく余接由来の増強 $g_{\text{obs}} = g_{\text{bar}} \cdot \eta(r)$ に従うはずであり、$g \lesssim a_0/10$ で測定可能なずれをもつ。
  \item \textbf{$R(M)$ からの多様性}：観測される回転曲線形状の多様性、すなわち緩やかに上昇するもの（$r_{\text{disk}} \ll R$ の低表面輝度矮小銀河）から平坦なもの（天の川型）から緩やかに減少するもの（$r_{\text{disk}} \sim R$ の高表面輝度渦巻）まで、は単一関係 $R = \sqrt{GM/a_0}$ によって追加の自由パラメーターなしに再現される。
\end{itemize}

暗黒物質は不要となる。コンパクトな双対時空 $S^3$ 幾何を、架空の平坦かつ無限の背景へと投影した産物である。余接ポテンシャル、方位正距歪み、および $S^3$ 境界におけるねじれ反転が、不可視物質なしにすべての不足重力を提供する。

\section{重力工学：双対時空ローターのコヒーレント励起}
\label{sec:gravitycontrol}

二重時空理論は、重力を侵すべからざる幾何的拘束から工学可能なねじれ自由度へと変換する。ねじれスカラー
\[
J = \frac{1}{16} B(\Omega_{\text{biv}}, \Omega_{\text{biv}}),\qquad
\Omega_{\text{biv}} = \log(R_{\text{usual}}^\dagger R_{\text{dual}})
\]
は、粒子内在的な通常ローターと双対ローターの間の相対的ずれのみに依存する。したがって双対ローターパラメーター $\theta_a$ をコヒーレントに変調できる外場は、$J \to 0$（重力遮蔽または慣性質量低減）あるいは $J < 0$（斥力重力）を駆動しうる。

\subsection{円偏光電磁場の双対ローターへの理論的結合}

電磁場は二つの時空にわたって自然に分割される。

\[
F = F_{\text{usual}} + F_{\text{dual}},
\]

ここで

\[
F_{\text{usual}} = \mathbf{E} \cdot (iI, iJ, iK) \quad (\hat{p}^2 = +1:\ \text{ブースト様、非コンパクト}),
\]

\[
F_{\text{dual}} = \mathbf{B} \cdot (I, J, K) \quad (\hat{q}^2 = -1:\ \text{回転様、コンパクト}).
\]

電場成分 $\mathbf{E}$ は主として通常時空ブースト生成子に結合し、磁場成分 $\mathbf{B}$ は専ら双対時空回転生成子に結合する。この代数的分離は、時間反転の下での相異なる変換性質（$T$: $\mathbf{E} \to \mathbf{E}$, $\mathbf{B} \to -\mathbf{B}$）の幾何的起源である。双対時空は内在的に時間反転している。

ヘリシティ $\sigma = \pm 1$ で $z$ に沿って伝播する円偏光波に対して、

\[
\mathbf{E} = E_0 (\hat{x} \cos(kz - \omega t) + \sigma \hat{y} \sin(kz - \omega t)), \quad \mathbf{B} = \frac{E_0}{c} (-\hat{x} \sin(kz - \omega t) + \sigma \hat{y} \cos(kz - \omega t)).
\]

時間平均は $\mathbf{E}$ の振動するブースト様寄与を消去し、純粋な双対時空回転励起を残す。

\[
\langle F_{\text{dual}} \rangle \propto \sigma E_0^2 \, K.
\]

ヘリシティ $\sigma$ は双対ローター駆動の方向を直接定める。右円偏光（$\sigma = +1$）は引力ねじれ不整合（$J > 0$）を強化し、左円偏光（$\sigma = -1$）はそれに対抗し、スカラー $J$ を低減または反転させる。このヘリシティ選択性は理論の最も鋭い現れである。重力は単に変調されるのではなく、双対時空の内在的時間反転非対称を通じてカイラルにアドレスされる。

結合は最小かつゲージ不変である。

\[
R_{\text{dual}} \to \mathcal{P} \exp\!\left( \frac{1}{2} \int (\theta_a \Gamma_a + e \lambda A_a \Gamma_a) \, ds \right),
\]

ここで双対生成子 $\Gamma_a$ が相互作用を媒介する。共鳴近傍では、双対位相は
\[
\dot{\theta} = \kappa \sigma E_0^2 \sin(\omega t - \theta + \delta_\sigma)
\]
として発展し、ねじれ不整合角のコヒーレント励起を与える。

\subsection{共鳴周波数と超伝導体の役割}

内在的双対時空スケールはプランク的である（$\sim 10^{43}$ Hz）。巨視的コヒーレンス（$N \sim 10^{26}$ 粒子）は有効周波数を $10^{12}$--$10^{15}$ Hz へと赤方偏移させる。超伝導体はジョセフソンプラズマ共鳴および集団位相ロックを通じてそれを劇的にさらに低減し、現行技術で到達可能な GHz--THz 範囲の可変モードを達成する。

\subsection{提案実験プロトコルと予測}

高 $Q$ 超伝導共振器内のマイクログラムスケール試験質量を、可変円偏光放射（0.1--10 THz）で照射する。期待されるシグネチャは、見かけの重量または慣性質量の共鳴的・ヘリシティ依存変化（$\Delta m/m \gtrsim 10^{-6}$）であり、左円偏光は低減を、右円偏光は増強を生じる。この水準での成功は実用的重力工学へと直接スケールする。

史上初めて、重力は単に制御可能であるだけでなくカイラルに制御可能となり、ヘリシティが二重時空のねじれ幾何における引力と斥力の間のスイッチとして働く。

\section{バリオン非対称：双対ローター時間矢固定からの幾何的起源}
\label{sec:baryonasymmetry}

二重時空理論において、宇宙の観測されるバリオン非対称、すなわちバリオンが反バリオンを圧倒的に上回る深い不均衡であり、パラメーター $\eta \equiv (n_B - n_{\bar{B}})/n_\gamma \approx 6 \times 10^{-10}$ によって定量化されるものは、プランク期における双対ローターの時間矢の非対称固定の不可避な帰結として現れる。この機構は双四元数代数 $\mathrm{Cl}(3,1)$ の内在的カイラリティに根ざし、付加的場、パラメーター、微調整を呼び出すことなくすべてのサハロフ条件を満たす。ここでは非対称を代数的に導き、通常時空と双対時空の間の幾何的相互作用を強調する。

\subsection{内在的時間デュアリティとローターパラメーター}

各粒子が16実次元双四元数構造の中にミンコフスキー時空の対を符号化することを想起する。未来向き時間基底 $j$ をもつ $\{j, kI, kJ, kK\}$ で張られる通常時空と、擬スカラー $i$ による右乗算で得られるその双対 $\{k, jI, jJ, jK\}$ であり、$j i = -k$ したがって過去向き時間基底 $k$ を与える。この双対操作 $X \mapsto X i$ はミンコフスキーノルムを保つ（第\ref{sec:math}節。明示的非可換計算が $(X i)^2 = X^2$ を与える）が時間的向きを反転させ、基本的カイラリティを刻印する。粒子（通常支配）は時間前方に伝播し、反粒子（双対支配）は動力学的抑制がなければ本質的に後方に伝播するであろう。

完全な運動学的発展は完全ローターに支配される。真空ではねじれなし（$J=0$）が要求され、時間矢の完全な反同期を保証する。前方の通常ローターは逆行する双対ローターによって映され、粒子--反粒子等価性を保つ。

\subsection{プランク期における時間矢の固定}

プランクスケール（$t \sim 10^{-43}$ s、$T \sim 10^{19}$ GeV）では、量子重力揺らぎ、すなわち双対ローターバイベクトルにおけるゼロ点振動として現れるものが、ローターパラメーターの真空期待値における相転移を通じて自発的対称性の破れを誘起する。双四元数基底の非対称、すなわち $j$ が唯一の時間様生成子であり空間基底が交差項（$kI$ など）を含むことが、一意のエネルギー的に好ましい配置を選択する。

\[
\langle \theta \rangle = \phi + \delta, \quad \delta > 0,
\]
ここで $\delta$ はプランク期に刻印された CP 奇のカイラル位相である。我々は $\delta$ を、その大きさが観測によって固定されるべき枠組みの小さな自由パラメーターとみなす。粗い見積もり $\delta \sim 10^{-5}$（観測される CKM 位相のヤールスコグ測度に匹敵）を通じて採用する。この固定は双対写像 $X \mapsto X i$ そのもののパリティ奇の性格を反映する。擬スカラー $i$ は空間反転の下で符号を反転する（$i \to -i$）ため、$R_\text{usual}$ と $R_\text{dual}$ を区別するいかなるプランクスケールポテンシャルも必然的に $\delta$ において CP 奇であり、$\delta$ の確定した符号をエネルギー的に好む。このポテンシャルの微視的形はここでは特定しない。それは本代数的議論の範囲外の量子重力動力学に支配され、単一のスカラーパラメーター $\delta > 0$ を通じてのみ入る。

得られるねじれ不整合を抽出するため、\S\ref{subsec:dual-rotors} で確立した双対変換 $\theta \mapsto \theta i$ を適用する。双対ローターの双対時空への物理的効果はブースト（$(\hat q\, i)^2 = +1$）であるため、過回転 $\delta$ は $\Omega_\text{biv}$ へ純ブースト方向過剰として寄与する。既存の平衡（$\theta = \phi$）は、キリング形式のヌル錐上で $i\Gamma_a$ および $\Gamma_a$ 係数が釣り合うバイベクトルを与える（$J^{(0)} = 0$、真空のねじれなしと整合）。カイラル摂動はこの釣り合いを非対称に破る。ブースト軸 $\hat p = \sum_a p_a\, i\Gamma_a$ に沿って $\delta_a = \delta\, p_a$ と置くと、線形化された不整合はブーストのみの形
\[
\Omega_\text{biv} \approx \sum_a \delta_a\, i\Gamma_a \in \mathfrak{so}(3,1)
\]
に還元され、\S\ref{sec:gravity} の弱場線形化と構造的に同一である。そこで確立されたキリング形式の値（$B(i\Gamma_a, i\Gamma_b) = +8\,\delta_{ab}$）は直接
\[
J = \frac{1}{16} B(\Omega_\text{biv}, \Omega_\text{biv}) = \frac{1}{2}\sum_a \delta_a^2 = \frac{\delta^2}{2} > 0
\]
を与える。これは代数的核心において粒子--反粒子対称性を破る。通常支配モード（物質）に対しては、ねじれ不整合は正（$J > 0$）であり引力重力束縛を生じる。双対支配モード（反物質）に対しては、通常／双対の役割割り当てが反転する。逆変換 $\phi \mapsto -\phi\, i$ はカイラル位相 $\delta$ をブースト方向ではなく回転方向へ写すため、不整合は $\Omega_\text{biv} \approx \sum_a \delta_a\, \Gamma_a$ となる。キリング形式の値 $B(\Gamma_a, \Gamma_b) = -8\,\delta_{ab}$ はしたがって $J = -\delta^2/2 < 0$ を与え、そのようなモードを宇宙論的真空から動力学的に排除する斥力ねじれ層である。

\subsection{サハロフ条件とオーダー評価}

固定はローター指数を通じて三つのサハロフ条件を自動的に満たす。

1. バリオン数の破れ：バリオン数 $B$ は双対ローターにおける擬スカラー位相として符号化される。具体的には $B \propto \arg(\det R_\text{dual})$。不整合 $\delta \neq 0$ は $X \mapsto X i i = -X$（二重時間反転、実効的にバリオン反転）のような非保存過程を可能にし、したがって通常支配モードと双対支配モードのレートは $O(\delta)$ で異なる。

2. C および CP の破れ：双対写像 $i$ は内在的にカイラルであり（パリティの下で奇、$i \mapsto -i$ が空間反転の下で）、$\delta > 0$ は観測される CKM 位相に匹敵する CP 奇位相 $\phi_\text{CP} = \arg(\exp(i \delta \Gamma_a)) \sim \delta$ を導入する。この幾何的 CP 破れは普遍的であり、湯川固有の調整なしにすべてのフェルミオン世代に浸透する。

3. 熱平衡からの離脱：インフレーション膨張は系を平衡から遠くへ追いやり、CP 奇位相 $\delta$ が、双対支配モードがその通常支配パートナーに対して対消滅するにつれて凍結された数密度非対称へと変換されることを許す。標準的凍結見積もりはカイラル位相について二次の非対称
\[
\eta \sim \delta^2
\]
を与え、したがって観測値 $\eta \sim 6 \times 10^{-10}$ は $\delta \sim 10^{-5}$ によって再現され、上記の CP 奇位相見積もりと一致する。精密な数値係数はインフレーション後の再加熱温度に依存し、この水準では予測されない。

インフレーション後、残余の双対支配反粒子はバリオン余剰に対して対消滅し、$\eta$ を時間矢固定の凍結遺物として残す。反粒子は「消える」のではなく動力学的に排除される。それらの後方時間矢は $\delta > 0$ によって固定された前方向き宇宙膨張と両立しない。

\subsection{観測的シグネチャと反証可能性}

この機構は、$\Omega_\text{biv}$ を源とするカイラル重力波からの微妙な CMB 異方性を予測し、テンソル対スカラー比 $r \sim \delta^2 \sim 10^{-10}$ および LiteBIRD のような将来ミッションで検出可能な CP 奇パリティスペクトルをもつ。低エネルギーでは、普遍的レプトン非対称 $\eta_\ell \approx \eta_B$ を含意し、ニュートリノレス二重ベータ崩壊実験（例：LEGEND-200）によって分解可能である。これらのシグネチャの非観測は理論を反証し、確認は双対時空を宇宙論の基礎的枠組みへと高めるであろう。

本質において、バリオン非対称は謎ではなく、宇宙の時間の矢の選択の幾何的残響である。物質には前方、反物質には禁止。双対ローターはいったん固定されると、我々の宇宙が粒子のみの聖域であると宣告する。

\section{離散時間とコンパクト化されたねじれスカラー}
\label{sec:discretetime}

時空連続体仮説の拒否は、基本スケールにおける離散構造を要求する。双四元数代数に符号化された内在的デュアリティ、とりわけ擬スカラー $i$ と、時間矢を反転させつつミンコフスキーノルムを保つ双対写像 $X \mapsto Xi$ は、双対時空を内在的にコンパクト化し、通常時空は非コンパクトのまま残す。

したがって双対時空回転角のみを離散化する。
\[
\theta \in \frac{2\pi}{N} \mathbb{Z},
\]
ここで $N \gg 1$ は双対時空のコンパクト化スケールを定める非常に大きな整数である（典型的にはプランク単位における粒子の静止質量に比例する）。通常時空ラピディティは連続のままである。
\[
\phi \in \mathbb{R}.
\]

この非対称は理論的整合性のために要求される。
\begin{itemize}
  \item 双対時空は楕円的（$\hat{q}^2 = -1$）であり本質的にコンパクトである。その離散化は $R_{\text{dual}}$ によってパラメトライズされる $S^3$ の $4\pi$ 周期位相を直接反映する。これはニュートン--ウィグナー局在化問題を解消し自然な量子離散性を与えるために本質的な有限位相体積を提供する。
  \item 通常時空は双曲的（$\hat{p}^2 = +1$）であり内在的に非コンパクトである。したがって $\phi$ の連続性は完全ローレンツ群を保ち、巨視的極限における古典一般相対性理論との正確な等価性を保証する。
\end{itemize}

ねじれ不整合ローター $\Omega = R_{\text{usual}}^\dagger R_{\text{dual}}$ はしたがって連続な通常成分と離散な双対成分を結合する。パラメータ空間スカラー $J=\frac12\sum_a(\alpha_a^2-\beta_a^2)$ は、許容配置上でのみ $|J|\le\Jbound$（同値に $|\Jnorm|\le 1$）で有界である。双対角のみの離散化は有界を含意しない。非負ラピディティなしの主枝条件は反例を許し、連続な通常セクターラピディティは非有界である。したがって連続体誘起の特異点が排除されるのは、許容錐が課された後に限られる。

ねじれ層にわたる有効ローレンツ因子において臨界的問題が生じる。それは典型的に
\[
\gamma_{\text{eff}} \propto D_N(\delta\theta) \equiv \frac{\sin(N \delta\theta / 2)}{N \sin(\delta\theta / 2)}
\]
の形の同期関数として現れ、ここで $\delta\theta$ はねじれ不整合角、$D_N$ は（正規化された）ディリクレ核である。分母は格子 $\delta\theta \in 2\pi \mathbb{Z}$ 上で正確に消える。

そのような各点で分子 $\sin(N \cdot k\pi)$ も消える（任意の整数 $N$、素数でも合成数でも）ため、見かけの特異点は実際にはディリクレ核型の除去可能な $0/0$ である。標準的なロピタル評価は有限極限
\[
\lim_{\delta\theta \to 2\pi k} D_N(\delta\theta) = \pm 1 \qquad (k \in \mathbb{Z})
\]
を与える。したがって同期関数はコンパクト格子全体で有界であり、古典連続体極限 $N \to \infty$ は病理なしに回復される。

したがってねじれ動力学の有界性はディリクレ核の構造的性質であり、$N$ が素数であることを要求しない。任意の大きな正整数 $N$ が、連続体極限を保ちつつ偽の特異点を抑制するのに十分である。双対時空コンパクト化の数論的構造、とりわけ $N$ と粒子質量スペクトルの関係は、枠組みの自然な拡張であり、離散双対時空模型に関する姉妹論文で詳細に展開される。

ディリクレ核の議論は \(N\) の素数性がねじれ動力学の有界性に厳密には要求されないことを示すが、双対時空コンパクト化のより深い算術構造は、\(N\) が奇素数（または十分大きな奇素数の制御された積）に選ばれるとき実質的に豊かになる。この場合、双四元数整数環は一意分解を継承し、許容されるねじれ層および質量スペクトルは素元によって生成されるディオファントス方程式の解となる。この「奇素数離散化」は、上記で確立された有界ねじれ動力学を保ちつつ、離散コンパクト化がよく定義された算術構造を許すことを保証する。

同一の離散--連続相補性はゼータ関数正則化を通じて精密化される。すべての素数にわたる正則化された無限積
\[
\prod_p p \;\overset{\text{zeta regularization}}{\longrightarrow}\; 4\pi^2
\]
はオイラー積の解析接続に起源をもち、素数の離散乗法構造が関数方程式に現れる連続因子 \(\pi\) および \(\Gamma\) によって補完される。二重時空枠組みでは、この数値因子は双対時空ローター位相を支配する基本角周期 \(2\pi\) の平方として幾何的に生じる。双対ローターは半角（\(\theta/2\)）を通じて定義され、三つの独立な回転平面がキリング形式を通じて寄与するため、有効周期性 \(2\pi\) は平方され、正確に \(4\pi^2\) を与える。これは素数積のゼータ正則化に現れる同一定数の直接的幾何実現を提供する。したがって \(N\) を奇素数に選ぶことは、離散--連続デュアリティが、さもなければ発散する量のよく定義され有限で物理的に意味のある正則化を生じることを保証する。

\section{双対ローター動力学とド・ブロイ位相の幾何的起源}
\label{sec:dualrotordynamics}

二重時空理論において、ねじれ不整合スカラー $J$ は相対ローター $\Omega = R_{\text{usual}}^\dagger R_{\text{dual}}$ から構成される。$J$ は古典場理論における重力作用密度として働く一方、個々の質量粒子の内部自由度を支配するポテンシャルとしての自然な粒子内在的解釈も許す。

小さな不整合角に対し、キリング形式不変量は
\[
J \approx \frac{1}{2} \sum_{a=1}^3 \delta\phi_a^2
\]
と展開される。ここで $\delta\phi_a = \phi_a - \theta_a$ は各ブースト--回転平面におけるねじれ不整合角である。静止質量 $m$ を双対セクターの剛性定数として取り込むと
\[
J = \frac{1}{2} m \sum_{a=1}^3 \delta\phi_a^2
\]
となる。

単一粒子の内在的ローター角に対する有効作用を提案する。
\[
S_{\text{particle}} = \int \left[ \frac{1}{2} \sum_{a=1}^3 \dot{\phi}_a^2 - \frac{1}{2} \sum_{a=1}^3 \dot{\theta}_a^2 - \frac{m}{2} \sum_{a=1}^3 (\phi_a - \theta_a)^2 \right] dt.
\]
運動項の反対符号は、通常セクターのブースト様性格と双対セクターの時間反転した回転様性格を反映する。

オイラー--ラグランジュ方程式は
\begin{align}
\ddot{\phi}_a &= m (\phi_a - \theta_a), \\
-\ddot{\theta}_a &= m (\phi_a - \theta_a).
\end{align}
ねじれ不整合角 $\delta\phi_a = \phi_a - \theta_a$ を定義すると、単純調和振動子方程式
\begin{equation}
\ddot{\delta\phi}_a + m \delta\phi_a = 0.
\label{eq:dual-harmonic}
\end{equation}
が得られる。振動周波数は $\sqrt{m}$ である（$\hbar = c = 1$ の自然単位）。

一様運動する自由粒子に対し、通常ラピディティ成分は一定率 $\dot{\phi}_a = \gamma v_a/c$ で発展する。式~\eqref{eq:dual-harmonic} の一般解は
\[
\delta\phi_a(t) = A_a \sin(\sqrt{m}\, t + \psi_a).
\]
である。低エネルギー非相対論的レジームでは、急速なコンプトン周波数振動は平均してゼロとなり、固有時に比例する位相の遅い線形蓄積を残す。$\hbar$ および $c$ を復元すると、蓄積位相は
\[
\int \delta\phi_a \, dt \propto \frac{E t - \mathbf{p} \cdot \mathbf{x}}{\hbar}
\]
となり、これは正確に物質波のド・ブロイ位相である。

したがって、量子力学的波動関数位相 $S(x,t)/\hbar$ は、通常ローターに対する双対ローター遅れの緩変化成分として幾何的に現れる。これは二重時空枠組み内で、付加的要請を呼び出すことなく、量子物質波の完全に背景独立な代数的起源を提供する。

この粒子内在的作用の多数粒子にわたる集団平均は古典遠平行作用 $\int J \, d^4x$ を回復し、巨視的スケールにおける一般相対性理論との完全な整合性を保証しつつ、双対時空自由度の微視的量子動力学を明らかにする。

この調和振動子動力学は、重力の古典ねじれ記述と次節で展開する完全量子理論との橋渡しを構成する。

\section{量子重力と完全統一に向けて}
\label{sec:quantumgravity}

これまでに提示した古典定式化は、重力が粒子内在的な通常時空と双対時空の間のねじれ不整合として現れ、$\mathrm{Cl}(3,1)$ と同型の16次元双四元数代数に符号化されることを示す。作用 $S$ は連続体仮説とリーマン曲率を捨てつつ、アインシュタイン--ヒルベルト作用との正確な等価性を与える。

深い拡張は、双対ローター $R_\text{dual}$ が数学的に SU(2) スピン $1/2$ 回転演算子と同一であることの認識から生じる。したがって双対セクター位相 $\theta$ は量子状態を直接パラメトライズする。双対時空の内在的コンパクト性（$\theta \mod 4\pi$）により、これらの位相はコンパクト多様体上に住み、各粒子の量子状態 $|\psi\rangle$ が $R_\text{dual}$ のユニタリ位相部分に対応するヒルベルト空間構造を自然に与える。

位置および運動量演算子は双対ローター角の射影として現れ、重力そのものを含む力は相対ローター位相の変調として現れる。これは量子力学を双対時空回転と、重力を同一代数枠組みにおけるねじれ不整合と同定し、有界な離散ねじれ層を通じて特異点を排除する。

明示的ヒルベルト空間構成、$\mathrm{Cl}(3,1)$ の極小左イデアルへのスピノル埋め込み、すべての相互作用の統一、およびプランク期時間矢固定を通じたバリオン非対称の解消を含む完全な量子重力枠組みは、姉妹論文「Quantum Gravity as Torsional Mismatch in Biquaternion Double Spacetime」\cite{dst-quantum-gravity} で詳細に展開される。

この拡張は理論を完成させる。一般相対性理論は古典的に保たれつつ、量子重力は余剰次元やアドホックな量子化なしに、すべての質量粒子に内在的な二重時空構造のみから代数的に現れる。

\section{結論}
\label{sec:conclusions}
双四元数代数 $\mathbb{H} \otimes_{\mathbb{R}} \mathbb{H} \cong \mathrm{Cl}(3,1)$ に根拠をもつ二重時空理論は、重力、慣性、基本相互作用を理解するための変革的枠組みを提供する。各粒子を絡み合ったミンコフスキー時空の対、すなわち通常と双対の中に埋め込むことにより、理論は重力現象をこれらのセクター間のねじれ不整合の現れとして再解釈する。このパラダイム転換は重力物理学における長年のパラドックスを解消し、強い核力を層状ねじれ動力学の現れとして統一し、内在的カイラリティと時間矢固定を通じたバリオン非対称の自然な機構を提供する。双対ローター同期を通じた重力制御から、原子核の原始ねじれ星としての再解釈に至る理論の予測は、実験的検証と技術革新の新たな道を開く。重力物理学の新時代の瀬戸際に立つ今、双対時空枠組みは宇宙と我々のその中での位置についての理解を再考するよう招く。

\bibliographystyle{unsrt}
\begin{thebibliography}{1}
\raggedright

\bibitem{einstein1928}
A.~Einstein,
``Riemann-Geometrie mit Aufrechterhaltung des Begriffes des Fernparallelismus,''
\textit{Sitzungsber. Preuss. Akad. Wiss. Phys.-Math. Kl.},
217--221 (1928).

\bibitem{einstein1939}
A.~Einstein,
``On a stationary system with spherical symmetry consisting of many gravitating masses,''
\textit{Ann. Math.},
\textbf{40}, no.~4, 922--936 (1939).
doi:10.2307/1968902

\bibitem{girard1999}
P.~R. Girard,
``Einstein’s equations and Clifford algebra,''
\textit{Advances in Applied Clifford Algebras},
\textbf{9}, no.~2, 225--230 (1999).
doi:10.1007/BF03042377

\bibitem{maluf2013}
J.~W. Maluf,
``The teleparallel equivalent of general relativity,''
\textit{Ann.\ Phys.\ (Berlin)} \textbf{525}, 339--359 (2013).
doi:10.1002/andp.201200272

\bibitem{dst-quantum-gravity}
Anonymous,
``Quantum Gravity as Torsional Mismatch in Biquaternion Double Spacetime'' (2025).
\url{https://github.com/hypernumbernet/dual-spacetime-doc/blob/main/dst-quantum-gravity.pdf}.

\end{thebibliography}

\clearpage
\appendix
\section{乗算表}
\label{app:mult-table}

\begin{table}[ht]
  \begin{center}
    \caption{双四元数乗算表}
    \small
    \begin{tabular}{r|rrrrrrrrrrrrrrr}
      &$j$&$kI$&$kJ$&$kK$&$iI$&$iJ$&$iK$&$I$&$J$&$K$&$k$&$jI$&$jJ$&$jK$&$i$ \\ \hline
      j&$-1$&$+iI$&$+iJ$&$+iK$&$-kI$&$-kJ$&$-kK$&$+jI$&$+jJ$&$+jK$&$+i$&$-I$&$-J$&$-K$&$-k$ \\ \hline
      kI&$-iI$&$+1$&$-K$&$+J$&$-j$&$+jK$&$-jJ$&$-k$&$+kK$&$-kJ$&$-I$&$+i$&$-iK$&$+iJ$&$+jI$ \\ \hline
      kJ&$-iJ$&$+K$&$+1$&$-I$&$-jK$&$-j$&$+jI$&$-kK$&$-k$&$+kI$&$-J$&$+iK$&$+i$&$-iI$&$+jJ$ \\ \hline
      kK&$-iK$&$-J$&$+I$&$+1$&$+jJ$&$-jI$&$-j$&$+kJ$&$-kI$&$-k$&$-K$&$-iJ$&$+iI$&$+i$&$+jK$ \\ \hline
      iI&$+kI$&$+j$&$-jK$&$+jJ$&$+1$&$-K$&$+J$&$-i$&$+iK$&$-iJ$&$-jI$&$-k$&$+kK$&$-kJ$&$-I$ \\ \hline
      iJ&$+kJ$&$+jK$&$+j$&$-jI$&$+K$&$+1$&$-I$&$-iK$&$-i$&$+iI$&$-jJ$&$-kK$&$-k$&$+kI$&$-J$ \\ \hline
      iK&$+kK$&$-jJ$&$+jI$&$+j$&$-J$&$+I$&$+1$&$+iJ$&$-iI$&$-i$&$-jK$&$+kJ$&$-kI$&$-k$&$-K$ \\ \hline
      I&$+jI$&$-k$&$+kK$&$-kJ$&$-i$&$+iK$&$-iJ$&$-1$&$+K$&$-J$&$+kI$&$-j$&$+jK$&$-jJ$&$+iI$ \\ \hline
      J&$+jJ$&$-kK$&$-k$&$+kI$&$-iK$&$-i$&$+iI$&$-K$&$-1$&$+I$&$+kJ$&$-jK$&$-j$&$+jI$&$+iJ$ \\ \hline
      K&$+jK$&$+kJ$&$-kI$&$-k$&$+iJ$&$-iI$&$-i$&$+J$&$-I$&$-1$&$+kK$&$+jJ$&$-jI$&$-j$&$+iK$ \\ \hline
      k&$-i$&$-I$&$-J$&$-K$&$+jI$&$+jJ$&$+jK$&$+kI$&$+kJ$&$+kK$&$-1$&$-iI$&$-iJ$&$-iK$&$+j$ \\ \hline
      jI&$-I$&$-i$&$+iK$&$-iJ$&$+k$&$-kK$&$+kJ$&$-j$&$+jK$&$-jJ$&$+iI$&$+1$&$-K$&$+J$&$-kI$ \\ \hline
      jJ&$-J$&$-iK$&$-i$&$+iI$&$+kK$&$+k$&$-kI$&$-jK$&$-j$&$+jI$&$+iJ$&$+K$&$+1$&$-I$&$-kJ$ \\ \hline
      jK&$-K$&$+iJ$&$-iI$&$-i$&$-kJ$&$+kI$&$+k$&$+jJ$&$-jI$&$-j$&$+iK$&$-J$&$+I$&$+1$&$-kK$ \\ \hline
      i&$+k$&$-jI$&$-jJ$&$-jK$&$-I$&$-J$&$-K$&$+iI$&$+iJ$&$+iK$&$-j$&$+kI$&$+kJ$&$+kK$&$-1$ \\ \hline
    \end{tabular}
  \end{center}
\end{table}
\begin{table}[ht]
  \begin{center}
    \caption{Cl(3,1) 変種乗算表（$e_0, e_1, e_2, e_3 \mapsto 0,1,2,3$）}
    \scriptsize
    \begin{tabular}{r|rrrrrrrrrrrrrrr}
      &$0$&$1$&$2$&$3$&$01$&$02$&$03$&$32$&$13$&$21$&$123$&$032$&$013$&$021$&$0123$\\\hline
      $0$&$-$&$+01$&$+02$&$+03$&$-1$&$-2$&$-3$&$+032$&$+013$&$+021$&$+0123$&$-32$&$-13$&$-21$&$-123$\\
      $1$&$-01$&$+$&$-21$&$+13$&$-0$&$+021$&$-013$&$-123$&$+3$&$-2$&$-32$&$+0123$&$-03$&$+02$&$+032$\\
      $2$&$-02$&$+21$&$+$&$-32$&$-021$&$-0$&$+032$&$-3$&$-123$&$+1$&$-13$&$+03$&$+0123$&$-01$&$+013$\\
      $3$&$-03$&$-13$&$+32$&$+$&$+013$&$-032$&$-0$&$+2$&$-1$&$-123$&$-21$&$-02$&$+01$&$+0123$&$+021$\\
      $01$&$+1$&$+0$&$-021$&$+013$&$+$&$-21$&$+13$&$-0123$&$+03$&$-02$&$-032$&$-123$&$+3$&$-2$&$-32$\\
      $02$&$+2$&$+021$&$+0$&$-032$&$+21$&$+$&$-32$&$-03$&$-0123$&$+01$&$-013$&$-3$&$-123$&$+1$&$-13$\\
      $03$&$+3$&$-013$&$+032$&$+0$&$-13$&$+32$&$+$&$+02$&$-01$&$-0123$&$-021$&$+2$&$-1$&$-123$&$-21$\\
      $32$&$+032$&$-123$&$+3$&$-2$&$-0123$&$+03$&$-02$&$-$&$+21$&$-13$&$+1$&$-0$&$+021$&$-013$&$+01$\\
      $13$&$+013$&$-3$&$-123$&$+1$&$-03$&$-0123$&$+01$&$-21$&$-$&$+32$&$+2$&$-021$&$-0$&$+032$&$+02$\\
      $21$&$+021$&$+2$&$-1$&$-123$&$+02$&$-01$&$-0123$&$+13$&$-32$&$-$&$+3$&$+013$&$-032$&$-0$&$+03$\\
      $123$&$-0123$&$-32$&$-13$&$-21$&$+032$&$+013$&$+021$&$+1$&$+2$&$+3$&$-$&$-01$&$-02$&$-03$&$+0$\\
      $032$&$-32$&$-0123$&$+03$&$-02$&$+123$&$-3$&$+2$&$-0$&$+021$&$-013$&$+01$&$+$&$-21$&$+13$&$-1$\\
      $013$&$-13$&$-03$&$-0123$&$+01$&$+3$&$+123$&$-1$&$-021$&$-0$&$+032$&$+02$&$+21$&$+$&$-32$&$-2$\\
      $021$&$-21$&$+02$&$-01$&$-0123$&$-2$&$+1$&$+123$&$+013$&$-032$&$-0$&$+03$&$-13$&$+32$&$+$&$-3$\\
      $0123$&$+123$&$-032$&$-013$&$-021$&$-32$&$-13$&$-21$&$+01$&$+02$&$+03$&$-0$&$+1$&$+2$&$+3$&$-$\\
    \end{tabular}
  \end{center}
\end{table}
これらは実際に同一の結果であり同型である。
\clearpage
\section{\texorpdfstring{ねじれ不整合スカラー $J$ と TEGR の遠平行ねじれスカラー $T$ の等価性}{ねじれ不整合スカラー J と TEGR の遠平行ねじれスカラー T の等価性}}
\label{app:TEGR-equivalence}

ねじれ不整合バイベクトルは
\[
\Omega_{\rm biv} = \log(R_{\rm usual}^\dagger R_{\rm dual}) = \sum_{a=1}^3 \left( \frac{\alpha_a}{2} i\Gamma_a + \frac{\beta_a}{2} \Gamma_a \right),
\]
であり、\( (i\Gamma_a)^2 = +1 \)、\( \Gamma_a^2 = -1 \) である。同軸生成子は可換であり \([i\Gamma_a,\Gamma_a]=0\)。異軸の交換子は必ずしも消えない（例：\([i\Gamma_1,\Gamma_2]\neq 0\)）。

六つの生成子 \(\{i\Gamma_a,\Gamma_a\}\) は \(\mathfrak{so}(3,1)\) の候補コピーを張る。12次元のラベル \(\mathfrak{so}(3,1)\oplus\mathfrak{so}(3,1)\) はこの基底に対して次元的に不正確である。並進と合わせて、自然な周囲代数はポアンカレ候補 \(\mathfrak{iso}(3,1)\cong\mathfrak{so}(3,1)\ltimes\mathbb{R}^{3,1}\) である。

キリング型対 \( B(i\Gamma_a, i\Gamma_a) = +8 \)、\( B(\Gamma_a, \Gamma_a) = -8 \) を上記の半角展開上で評価すると
\[
B(\Omega_{\rm biv}, \Omega_{\rm biv}) = 2 \sum_{a=1}^3 (\alpha_a^2 - \beta_a^2)
\]
が得られ、本付録の初期草稿に書かれた因子 \(8\sum\) ではない。離散姉妹論文および Lean で用いられるパラメータ空間ねじれスカラーは
\[
J = \frac12 \sum_{a=1}^3 (\alpha_a^2 - \beta_a^2)
\]
であり、この展開の下で \(\frac1{16}B(\Omega_{\rm biv},\Omega_{\rm biv})\) の4倍である。有界性の主張はこの \(J\) および \(\Jnorm=\frac8{3\pi^2}J\) に関するものである。許容配置（非負ラピディティかつ \(\alpha_a+\beta_a\le\pi/2\)）においては \(|J|\le\Jbound\) かつ \(|\Jnorm|\le 1\) である。主枝条件のみを課した場合、この有界は成り立たない。

この不変量を TEGR と比較するには、まずローター場とテトラッドの間の幾何対応を指定しなければならない。具体的には、完全ローター場 $R_{\rm total}(x) = R_{\rm usual}(x)\, R_{\rm dual}(x)$ は固定参照テトラッド $\{\mathring{e}^a\}$ にスピノルサンドイッチ積
\[
e^a{}_\mu(x) \;=\; \langle R_{\rm total}(x)\, \mathring{e}^a\, R_{\rm total}^\dagger(x) \rangle_{1,\mu}
\]
によって作用する。ここで $\langle \cdot \rangle_{1,\mu}$ は座標基底に沿う1次（ベクトル）成分を抽出する。この同定は $R_{\rm total}$ を計量 $g_{\mu\nu} = \eta_{ab}\, e^a{}_\mu e^b{}_\nu$ をもつテトラッド場 $e^a{}_\mu$ へと昇格させ、双対不整合バイベクトル $\Omega_{\rm biv}$ を $e^a{}_\mu$ に付随するヴァイツェンベック接続のコントルションに比例させる。
\[
\Omega_{\rm biv}\big|_{\mu\nu}^{ab} \;\propto\; K^{ab}{}_\mu\, dx^\mu \wedge dx^\nu,
\qquad
T^a{}_{\mu\nu} \;=\; \partial_\mu e^a{}_\nu - \partial_\nu e^a{}_\mu.
\]
この同定の下で、二次不変量 $J$ は（全発散を除いて）遠平行ねじれスカラー
\[
T = \frac14 T_{\rho\mu\nu}T^{\rho\mu\nu} + \frac12 T_{\rho\mu\nu}T^{\nu\mu\rho} - T_\rho T^\rho
\]
に
\[
J = \frac12 T + \nabla_\mu B^\mu
\]
の形で対応すると予想される。この恒等式はここでは再導出しない。チャート上のシュヴァルツシルト切片のみが姉妹 PGA 論文で形式化されており、変分的 TEGR／アインシュタイン--ヒルベルト等価性は文献から取る。

同一の基礎テトラッドに対し、TEGR の標準的結果 \cite{maluf2013} として
\[
\int T\, e\, d^4x = -\int R\sqrt{-g}\, d^4x + \int \partial_\mu(\sqrt{-g}\, V^\mu)\, d^4x
\]
が成り立つ。すなわち $T$ と $-R$ は全発散によってのみ異なる。$J$ に対する関係を代入すると
\[
S = \frac{c^4}{16\pi G} \int J \, d^4x = -\frac{c^4}{16\pi G} \int R \sqrt{-g}\, d^4x + (\text{境界項})
\]
が得られる。

$S$ の変分は、テトラッド $e^a{}_\mu$（あるいは同値にローター場 $R_{\rm total}$）を独立な動力学変数として取り、すべての場の変分 $\delta e^a{}_\mu$ が積分領域の境界 $\partial\mathcal{M}$ 上で消えるという標準的境界条件のもとで行う。この条件の下で上記の全発散項は変分方程式に寄与せず、残りの変分はアインシュタイン方程式
\[
G_{\mu\nu} = \frac{8\pi G}{c^4} T_{\mu\nu}
\]
を再現する。

したがって、上記で述べたローターからテトラッドへの同定と標準的変分境界条件のもとで、二重時空作用は古典極限においてアインシュタイン--ヒルベルト作用と動学的に等価である。

\end{document}




